%% toc: P3 37 37.1
\chapter{Abstract Logics, Principal Examples, and Comparison Maps}
\label{sec:p03-w03-p03-047-abstract-logics}

A \emph{logic} in the sense of this chapter is not a proof calculus.  It is
an assignment of sentences to vocabularies together with a satisfaction
relation between structures and those sentences.  The purpose of the
chapter is to make that notion exact, to fix the closure conditions under
which logics can be compared, and to prove the elementary facts that relate
the familiar properties of first-order logic (compactness, the
L\"owenheim--Skolem property, interpolation, definability, effective
enumerability of validities) once those properties are stated for an
arbitrary logic.  We write \(\mathscr L\) for such a logic and
\(\models_{\mathscr L}\) for its satisfaction relation.

\emph{Inputs.} Work in \(\mathsf{ZFC}\).  Structures have nonempty domains.
Vocabularies are sets of sort, relation, function and constant symbols, and
the supply of symbols is unbounded: for every set of symbols there are
symbols of every kind and arity outside it.  First-order logic
\(\mathscr L_{\omega\omega}\) and its compactness and L\"owenheim--Skolem
theorems are taken from \cref{thm:p02-w02-p02-005-bw-compactness,thm:p02-w02-p02-005-bw-ls,thm:p02-w02-p02-033-upward-ls-inclusion}.
The non-finite-generation proof uses the standard-model graph results of
\cref{thm:p02-w02b-p02-034-pr-graph,thm:p05-rep-numeralwise-relations}
for its explicitly described local syntax coding.

\emph{Conclusions.} The chapter proves: the recursion-theoretic form of the
completeness problem (\cref{thm:p03-w03-p03-047-abstract-completeness}); that
a complete calculus with finite derivations yields compactness
(\cref{thm:p03-w03-p03-047-completeness-compactness}); the undefinability
facts behind Mostowski's cardinality quantifiers
(\cref{thm:p03-w03-p03-047-fo-undefinable}); the definitions of finiteness and
of the natural numbers in \(\mathscr L(Q_0)\) and in
\(\mathscr L_{\omega_1\omega}\)
(\cref{thm:p03-w03-p03-047-q0-omega,thm:p03-w03-p03-047-lomega1-q0}); the
existence of Hanf numbers (\cref{thm:p03-w03-p03-047-hanf-exists}); the
separation form of interpolation and \(\Delta\)-closure
(\cref{thm:p03-w03-p03-047-craig-separation}); boundedness of every
countably compact regular logic (\cref{thm:p03-w03-p03-047-bounded-omega});
and, for Ebbinghaus's general framework, the failure of the renaming
property for \(\omega\)-logic, elimination of function symbols, renaming
from substitution, and the Beth property from interpolation
(\cref{thm:p03-w03-p03-047-omega-logic-renaming,thm:p03-w03-p03-047-function-elimination,thm:p03-w03-p03-047-substitution-renaming,thm:p03-w03-p03-047-craig-beth}).
For the principal logics it proves that \(\mathscr L(Q_0)\) has no
completeness theorem (\cref{thm:p03-w03-p03-047-q0-not-re}), what
\(\mathscr L(Q_1)\) does and does not express
(\cref{thm:p03-w03-p03-047-q1-mostowski-facts}), a second-order sentence
valid exactly under the continuum hypothesis
(\cref{thm:p03-w03-p03-047-sol-ch}), the basic properties of weak
second-order logic (\cref{thm:p03-w03-p03-047-weak-sol}), downward
L\"owenheim--Skolem to \(\aleph_1\) for \(\mathscr L(Q_1)\)
(\cref{thm:p03-w03-p03-047-q1-ls-consequence}), interpolants for Keisler's
entailment in stronger logics (\cref{thm:p03-w03-p03-047-q1-interpolants}),
and the behaviour of quantification over countable and uncountable subsets
(\cref{thm:p03-w03-p03-047-set-quantifier-failures,thm:p03-w03-p03-047-pos-monotone}).

The general compactness criterion is \cref{thm:p03-general-pinning}.
\Cref{sec:p03-compactness-vocabulary} derives the essential-vocabulary
bound, proves the theory-space correspondence, and gives an explicit
counterexample to its naive fixed-cardinal version.
The class-quantifier development in \cref{sec:p03-lindstrom-quantifiers}
proves representation, domain translation and its closure boundaries,
finite and infinitary comparison theorems, and the impossibility of
finitely generating full second-order logic by these quantifiers.
It also supplies Keisler's non-interpolation comparison and distinguishes
standard-part semantics from adding a quantifier, including the finite
model boundary for same-domain projective reductions.

\emph{Sources.} The chapter follows two chapters of \emph{Model-Theoretic
Logics} in the 2017 Cambridge reissue (Perspectives in Logic~8) of the 1985
volume: Barwise's background chapter \cite[Ch.~I, pp.~3--23]{Barwise1985BackgroundAims}
and Ebbinghaus's general framework
\cite[Ch.~II, pp.~25--72]{Ebbinghaus1985ExtendedLogics}.
Page numbers are those printed in that reissue.  The regular relational
package used for Lindstr\"om's theorem is fixed in
\cref{sec:p03-lindstrom}; Ebbinghaus's more general many-sorted notion is
developed here and related to it in
\cref{thm:p03-w03-p03-047-lindstrom-bridge}.

\section{Logics that embody mathematical concepts}
\label{sec:p03-w03-p03-047-barwise-extended}

Barwise distinguishes two parts of the subject.  \emph{Extended model
theory} isolates particular logics built to study particular mathematical
notions; \emph{abstract model theory} studies relations between logics
\cite[Ch.~I, p.~3]{Barwise1985BackgroundAims}.  This section treats the first.

\begin{hbdefinition}{def:p03-w03-p03-047-barwise-logic}{A logic in Barwise's sense}
A \emph{logic} has three ingredients: some class of mathematical
structures, some class of formal expressions, and a satisfaction relation
\(\mathfrak M\models\varphi\) linking structures to expressions \cite[Ch.~I, \S1.1, p.~4]{Barwise1985BackgroundAims}.  A sentence of
the logic is \emph{valid} when it holds in every structure of the logic.
The interest lies in logics whose structures have some mathematical property
built in and whose expressions formalize talk about that property.
Ebbinghaus's \cref{def:p03-w03-p03-047-ebbinghaus-logic} is the exact
version used for every theorem below; Barwise's description is the informal
reading of it.
\end{hbdefinition}

\begin{hbdefinition}{def:p03-w03-p03-047-first-order-thesis}{The first-order thesis}
Barwise names the \emph{first-order thesis} as the view that logic is the
logic implicit in the ``logical constants'' (connectives, the two
quantifiers, and on some versions identity), so that whatever cannot be
defined in first-order logic lies outside logic
\cite[Ch.~I, \S1.2, p.~5]{Barwise1985BackgroundAims}.  It is a philosophical
position, not a mathematical statement; no theorem below assumes or refutes
it.  Its mathematical content for this chapter is the list of notions that
first-order logic cannot define (\cref{thm:p03-w03-p03-047-fo-undefinable}),
which is what extended model theory sets out to supply.
\end{hbdefinition}

\subsection{The completeness problem and compactness}

\begin{hbdefinition}{def:p03-w03-p03-047-abstract-completeness}{Abstract completeness}
Let a logic have countably many sentences in a fixed vocabulary, coded by
natural numbers so that the set of codes of sentences is decidable.  The
\emph{completeness problem} asks for an effective list of valid axioms and
of valid rules of inference that generate exactly the valid sentences.  The
logic is \emph{abstractly complete} when its set of valid sentences is
recursively enumerable \cite[Ch.~I, \S1.3, pp.~6--7]{Barwise1985BackgroundAims}.
Here an \emph{effective calculus} is a recursively enumerable set of axioms
together with a recursively enumerable set of rule instances, each instance
having finitely many premisses and one conclusion; a derivation is a finite
sequence each of whose members is an axiom or the conclusion of a rule
instance whose premisses occur earlier.
\end{hbdefinition}

\begin{hbtheorem}{thm:p03-w03-p03-047-abstract-completeness}{proved=annex:p03-w03-p03-047-abstract-completeness}{The recursion-theoretic form of the completeness problem}
For a logic coded as in \cref{def:p03-w03-p03-047-abstract-completeness}:
\begin{enumerate}
 \item if some effective calculus derives exactly the valid sentences, the
       set of valid sentences is recursively enumerable;
 \item if the set of valid sentences is recursively enumerable, then the
       calculus whose axioms are the valid sentences and which has no rules
       derives exactly the valid sentences.
\end{enumerate}
\end{hbtheorem}

The second clause is why Barwise says the converse gives a completeness
theorem only ``in principle'': the calculus it produces has no simple axioms
and rules, and finding a perspicuous one is a separate task
\cite[Ch.~I, \S1.3, p.~7]{Barwise1985BackgroundAims}.  A logic may lack every
such calculus: for full second-order logic the valid sentences are not
recursively enumerable (\cref{thm:p02-sol-full-no-effective-completeness}).

\begin{hbdefinition}{def:p03-w03-p03-047-compactness-properties}{Strong and countable compactness}
A logic has the \emph{strong compactness property} when every set \(T\) of
its sentences, each finite subset of which has a model, has a model.  It has
the \emph{countable compactness property} when the same holds for countable
sets \(T\) \cite[Ch.~I, \S1.4, p.~7]{Barwise1985BackgroundAims}.  For logics with
infinitely long proofs Barwise notes that ``finite'' is replaced by another
notion of small set, first by cardinality and later by notions from
generalized recursion theory \cite[Ch.~I, \S1.4, p.~8]{Barwise1985BackgroundAims};
the cardinal version is \(\kappa\)-compactness in
\cref{def:p03-w03-p03-047-ebbinghaus-properties}.
\end{hbdefinition}

\begin{hbtheorem}{thm:p03-w03-p03-047-completeness-compactness}{proved=annex:p03-w03-p03-047-completeness-compactness}{A complete calculus with finite derivations yields compactness}
Suppose a logic has a sentence \(\bot\) true in no structure, and a derivation
relation \(\vdash\) between sets of sentences and sentences such that
\begin{enumerate}
 \item \emph{soundness:} \(T_0\vdash\varphi\) implies that every model of
       \(T_0\) is a model of \(\varphi\);
 \item \emph{finite character:} if \(T\vdash\varphi\) then
       \(T_0\vdash\varphi\) for some finite \(T_0\subseteq T\);
 \item \emph{strong completeness:} if every model of \(T\) is a model of
       \(\varphi\), then \(T\vdash\varphi\).
\end{enumerate}
Then the logic has the strong compactness property.  If completeness is
assumed only for countable \(T\), the logic has the countable compactness
property.
\end{hbtheorem}

\subsection{Mostowski's quantifiers}

\begin{hbdefinition}{def:p03-w03-p03-047-mostowski-q}{The logics \(\mathscr L(Q_\alpha)\)}
Add to first-order syntax the formation rule: if \(\varphi(x)\) is a formula
then so is \(Qx\,\varphi(x)\), with \(x\) not free in the new formula; the
rule iterates with the connectives and the two quantifiers.  For an ordinal
\(\alpha\), the logic \(\mathscr L(Q_\alpha)\) gives \(Q\) the meaning
\[
 \mathfrak M\models Qx\,\varphi(x)
 \quad\text{iff}\quad
 \text{there are at least }\aleph_\alpha\text{ elements }b
 \text{ with }\mathfrak M\models\varphi(b)
\]
\cite[Ch.~I, \S1.5, p.~8]{Barwise1985BackgroundAims}, the proposal going back to
\cite{Mostowski1957}.  The logics \(\mathscr L(Q_\alpha)\) share one syntax
and differ only in the semantics of \(Q\).  Thus \(Q_0x\,\varphi\) says that
infinitely many elements satisfy \(\varphi\), and \(Q_1x\,\varphi\) that
uncountably many do (\cref{def:p03-q-one}).  Ebbinghaus writes the same
logics \(\mathscr L_{\omega\omega}(Q_\alpha)\)
\cite[Ch.~II, \S2.2, p.~34]{Ebbinghaus1985ExtendedLogics}.
\end{hbdefinition}

\begin{hbtheorem}{thm:p03-w03-p03-047-fo-undefinable}{proved=annex:p03-w03-p03-047-fo-undefinable}{What first-order logic cannot define}
Let \(P\) be a unary relation symbol.
\begin{enumerate}
 \item No first-order sentence \(\psi\) of vocabulary \(\{P\}\) satisfies,
       for all \(\{P\}\)-structures \(\mathfrak A\):
       \(\mathfrak A\models\psi\) iff \(P^{\mathfrak A}\) is infinite.  The
       same holds with ``finite'' in place of ``infinite''.  Only the
       countable compactness of first-order logic is used.
 \item No first-order sentence \(\psi\) of vocabulary \(\{P\}\) satisfies:
       \(\mathfrak A\models\psi\) iff \(P^{\mathfrak A}\) is uncountable.
       Only the downward L\"owenheim--Skolem theorem is used.
\end{enumerate}
\end{hbtheorem}

The corresponding failures for the extended logics, that \(\mathscr L(Q_0)\)
and \(\mathscr L_{\omega_1\omega}\) lack countable compactness and that
\(Q_1x\,(x=x)\) has a model but no countable one, are parts~3 and~4 of
\cref{thm:p02-w02-p02-005-bw-extended-ls}; the same theorem gives the
L\"owenheim property of \(\mathscr L_{\omega_1\omega}\) and of
\(\mathscr L(Q_0)\) for countable sets of sentences (parts~1 and~2).

These are the two observations from which Barwise draws Lindstr\"om's
theorem: each genuine strengthening of first-order logic loses the
first-order theorem that showed the new notion undefinable
\cite[Ch.~I, \S2.1, p.~14]{Barwise1985BackgroundAims}.

\begin{hbtheorem}{thm:p03-w03-p03-047-q0-omega}{proved=annex:p03-w03-p03-047-q0-omega}{\(\mathscr L(Q_0)\) characterizes the natural numbers}
Let \(\theta\) be the \(\mathscr L(Q_0)\)-sentence of vocabulary \(\{<\}\)
\[
 \text{``}<\text{ is a strict linear order''}
 \;\land\; Q_0x\,(x=x)
 \;\land\;\forall y\,\neg Q_0x\,(x<y).
\]
A \(\{<\}\)-structure satisfies \(\theta\) iff it is isomorphic to
\((\omega,<)\).
\end{hbtheorem}

This is one precise instance of Barwise's remark that \(\mathscr L(Q_0)\)
builds in the finite/infinite distinction and can define the natural numbers
\cite[Ch.~I, \S1.5, p.~8]{Barwise1985BackgroundAims}.  Since \(\theta\) is a
single sentence with a model of power \(\aleph_0\) and none of power
\(\aleph_1\), the upward L\"owenheim--Skolem theorem also fails for
\(\mathscr L(Q_0)\).

\begin{hbexample}{def:p03-w03-p03-047-q0-example}{Checking \(\theta\) on three orders}
In \((\omega,<)\) each element \(n\) has the \(n\) predecessors
\(0,\dots,n-1\), finitely many, and the domain is infinite; \(\theta\)
holds.  In \((\omega+1,<)\) the top element \(\omega\) has infinitely many
predecessors, so the clause \(\forall y\,\neg Q_0x\,(x<y)\) fails.  In the
finite order \((5,<)\) every element has finitely many predecessors but
\(Q_0x\,(x=x)\) fails.  The proof of
\cref{thm:p03-w03-p03-047-q0-omega} shows that these are the only two ways
to fail for a linear order.
\end{hbexample}

Barwise records further facts about these logics that are proved elsewhere.
Vaught's two-cardinal theorem gives an abstract completeness theorem for
\(\mathscr L(Q_1)\); Keisler found an explicit complete set of axioms and
rules; Fuhrken showed that \(\mathscr L(Q_1)\) is countably compact while
\(\mathscr L(Q_0)\) is not; and \(\mathscr L(Q_0)\) has no completeness
theorem \cite[Ch.~I, \S1.5, pp.~8--9]{Barwise1985BackgroundAims}.  Keisler's
completeness theorem, with his axioms and a full proof, is
\cref{thm:p02-w02c-p02-003-keisler-completeness}, and the countable
compactness of \(\mathscr L(Q_1)\) in countable vocabularies is its
\cref{thm:p02-w02c-p02-003-keisler-compactness}; Barwise notes that
Fuhrken's compactness result follows at once from Keisler's theorem, which
is the route taken there.  The transfer behind Vaught's theorem is discussed
after \cref{thm:p03-w03-p03-029-theory-form}; the recursive enumerability of
\(\mathscr L(Q_1)\)-validity is derived below from Keisler's calculus rather
than from the two-cardinal theorem
(\cref{thm:p03-w03-p03-047-q1-mostowski-facts}).  Barwise also describes
Lindstr\"om's general quantifiers \(Q_K\), one for each class \(K\) of
structures, and Keisler's probability quantifiers \((Px>r)\varphi(x)\),
whose logic is incomparable with first-order logic because measurability of
definable sets rules out closure under \(\forall\) and \(\exists\)
\cite[Ch.~I, \S1.5, p.~9]{Barwise1985BackgroundAims}.

\begin{hbtheorem}{thm:p03-w03-p03-047-q0-not-re}{proved=annex:p03-w03-p03-047-q0-not-re}{\(\mathscr L(Q_0)\) has no completeness theorem}
Let \(\Lang_A=\{0,S,+,\times\}\) and \(\tau_<=\Lang_A\cup\{<\}\), and let
\(\chi_{\mathcal N}\) be the \(\mathscr L(Q_0)\)-sentence of \(\tau_<\) that
is the conjunction of \(\theta\) of \cref{thm:p03-w03-p03-047-q0-omega} with
\[
 \forall x\,\neg\,x<0,\quad
 \forall x\forall y\,\bigl(Sx=y\leftrightarrow(x<y\land\neg\exists z\,(x<z\land z<y))\bigr),
\]
\[
 \forall x\,(x+0=x),\quad \forall x\forall y\,(x+Sy=S(x+y)),\quad
 \forall x\,(x\times0=0),\quad \forall x\forall y\,(x\times Sy=(x\times y)+x).
\]
\begin{enumerate}
 \item A \(\tau_<\)-structure satisfies \(\chi_{\mathcal N}\) iff it is
       isomorphic to \((\mathcal N,<)\), the standard model of arithmetic
       with its order.
 \item For every first-order \(\Lang_A\)-sentence \(\sigma\):
       \(\mathcal N\models\sigma\) iff \(\chi_{\mathcal N}\to\sigma\) is valid.
 \item Fix a coding of the \(\mathscr L(Q_0)\)-sentences of \(\tau_<\) by
       natural numbers such that the map sending the code of a first-order
       \(\Lang_A\)-sentence \(\sigma\), in the coding of
       \cref{thm:p02-w02-p02-040-no-effective-calculus}, to the code of
       \(\chi_{\mathcal N}\to\sigma\) is computable.  Then the set of valid
       \(\mathscr L(Q_0)\)-sentences of \(\tau_<\) is not recursively
       enumerable, and by \cref{thm:p03-w03-p03-047-abstract-completeness}
       no effective calculus derives exactly these valid sentences.
\end{enumerate}
\end{hbtheorem}

\begin{hbtheorem}{thm:p03-w03-p03-047-q1-mostowski-facts}{proved=annex:p03-w03-p03-047-q1-mostowski-facts}{What \(\mathscr L(Q_1)\) does and does not express}
\begin{enumerate}
 \item For a countable vocabulary in which the axiom schemes of
       \cref{def:p02-w02c-p02-003-kq-calculus} form a decidable set of codes,
       the valid \(\mathscr L(Q_1)\)-sentences form a recursively enumerable
       set.
 \item No \(\mathscr L(Q_1)\)-sentence \(\psi\) of vocabulary \(\{<\}\) has
       exactly the structures isomorphic to \((\omega,<)\) as models.  Hence
       \(\theta\) of \cref{thm:p03-w03-p03-047-q0-omega} has no equivalent
       in \(\mathscr L(Q_1)\), and \(\mathscr L(Q_1)\) does not include
       \(\mathscr L(Q_0)\).
 \item In the vocabulary \(\{\cdot,{}^{-1},e\}\) of groups, a group
       satisfies \(\neg Q_1x\,(x=x)\) iff it is countably generated, and
       satisfies \(Q_1x\,(x=x)\) iff it is uncountable.
\end{enumerate}
\end{hbtheorem}

Barwise names finite-dimensional vector spaces, torsion groups and
finitely generated groups among the notions that \(\mathscr L(Q_0)\) can
define, without fixing a vocabulary or giving a defining sentence
\cite[Ch.~I, \S1.5, p.~8]{Barwise1985BackgroundAims}.  For vector spaces the
claim depends on the field.  Write \(\tau_F=\{+,-,0\}\cup\{\lambda_a:a\in F\}\)
for the one-sorted vocabulary of vector spaces over a field \(F\), with
\(\lambda_a\) a unary function symbol for multiplication by the scalar
\(a\); a vector space over \(F\) is a \(\tau_F\)-structure in the usual way.

\begin{hbtheorem}{thm:p03-w03-p03-047-q0-vector-spaces}{proved=annex:p03-w03-p03-047-q0-vector-spaces}{Finite dimension in \(\mathscr L(Q_0)\)}
\begin{enumerate}
 \item If \(F\) is finite, a vector space over \(F\) is finite-dimensional
       iff it satisfies \(\neg Q_0x\,(x=x)\).
 \item Let \(T\) be the first-order \(\tau_{\mathbb Q}\)-theory of nonzero
       vector spaces over \(\mathbb Q\).  For every
       \(\mathscr L(Q_0)\)-formula \(\chi\) of \(\tau_{\mathbb Q}\) there is
       a first-order formula \(\chi^\circ\) of \(\tau_{\mathbb Q}\) with the
       same free variables such that \(\chi\) and \(\chi^\circ\) are
       satisfied by the same assignments in every model of \(T\).
 \item \(T\) is complete.  Hence \(\mathbb Q\) (dimension \(1\)) and
       \(\mathbb Q^{(\omega)}\) (countably infinite dimension) satisfy the
       same \(\mathscr L(Q_0)\)-sentences of \(\tau_{\mathbb Q}\), and no
       \(\mathscr L(Q_0)\)-sentence of \(\tau_{\mathbb Q}\) is true in
       exactly the finite-dimensional vector spaces over \(\mathbb Q\).
\end{enumerate}
\end{hbtheorem}

So the vector-space example holds over finite fields, where finite
dimension and finite cardinality coincide, and fails over \(\mathbb Q\) in
the one-sorted vocabulary: there \(Q_0\) adds nothing to first-order logic,
because every definable set of vectors is finite or cofinite.  Weak
second-order logic, which can quantify over finite spanning sets, is the
logic in which finite dimension is a natural notion; the torsion-group
definition in that logic is part~2 of \cref{thm:p03-w03-p03-047-weak-sol}.
For torsion groups and finitely generated groups in the vocabulary of
groups, no \(\mathscr L(Q_0)\)-definition and no proof of undefinability is
given here.  The natural numbers are \cref{thm:p03-w03-p03-047-q0-omega},
and the arithmetic structure \cref{thm:p03-w03-p03-047-q0-not-re}.

\subsection{Infinitary and second-order logic}

\begin{hbdefinition}{def:p03-w03-p03-047-lkl}{The infinitary logics \(\mathscr L_{\kappa\lambda}\)}
For infinite cardinals \(\kappa\) and \(\lambda\), \(\mathscr L_{\kappa\lambda}\)
is generated from atomic formulas by negation, conjunctions and disjunctions
of sets of fewer than \(\kappa\) formulas, and homogeneous strings of fewer
than \(\lambda\) quantifiers \cite[Ch.~I, \S1.6, p.~10]{Barwise1985BackgroundAims}.
The case \(\mathscr L_{\omega_1\omega}\) allows countable conjunctions and
disjunctions and finite quantifier strings; its syntax and satisfaction are
\cref{def:p02-lomega-one-omega}.
\end{hbdefinition}

\begin{hbtheorem}{thm:p03-w03-p03-047-lomega1-q0}{proved=annex:p03-w03-p03-047-lomega1-q0}{\(Q_0\) is definable in \(\mathscr L_{\omega_1\omega}\)}
For every \(\mathscr L_{\omega_1\omega}\)-formula \(\varphi(y)\) (possibly
with further free variables, which are held fixed), the formula
\[
 \bigwedge_{n\ge 0}\forall x_1\ldots x_n\,\exists y\,
 (\varphi(y)\land y\ne x_1\land\cdots\land y\ne x_n)
\]
holds in a structure iff infinitely many elements satisfy \(\varphi\)
\cite[Ch.~I, \S1.6, pp.~10--11]{Barwise1985BackgroundAims}.  Consequently every
\(\mathscr L(Q_0)\)-formula has an equivalent \(\mathscr L_{\omega_1\omega}\)-formula.
\end{hbtheorem}

The case \(n=0\) of the conjunction is \(\exists y\,\varphi(y)\).  Barwise
reports the further history: Hanf showed that a compact
\(\mathscr L_{\kappa\lambda}\) in the hoped-for sense requires strong
set-theoretic assumptions; Karp gave a completeness theorem for
\(\mathscr L_{\omega_1\omega}\) (\cref{thm:p02-karp-completeness});
Lopez-Escobar proved interpolation and a Beth theorem for it from a
Gentzen-style completeness proof (the two theorems are
\cref{thm:p02-w02c-p02-003-interpolation,thm:p02-w02c-p02-003-beth});
Mostowski showed that a logic with finite
syntax in which finiteness is definable has implicitly definable notions that
are not explicitly definable; and Moschovakis made every inductive definition
explicit by an infinite alternating string of quantifiers
\[
 \forall x_1\exists y_1\ldots\forall x_n\exists y_n\ldots
 \bigwedge_n\varphi_n(x_1,y_1,\ldots,x_n,y_n),
\]
generalizing a theorem of Svenonius on PC-classes of countable models
\cite[Ch.~I, \S1.6, pp.~10--11]{Barwise1985BackgroundAims}.  The moral Barwise
draws from Mostowski's result is that a logic which expresses finiteness and
is closed under implicit definability needs infinitary syntax in some sense
\cite[Ch.~I, \S1.6, p.~10]{Barwise1985BackgroundAims}.

For \(\mathscr L(Q_0)\) itself the failure of the Beth property can be
exhibited with the characterization of \((\omega,<)\) and nothing more.
Write \(S(x,y)\) for the first-order formula
\(x<y\land\neg\exists z\,(x<z\land z<y)\) (``\(y\) is the immediate successor
of \(x\)'').

\begin{hbtheorem}{thm:p03-w03-p03-047-q0-beth-failure}{proved=annex:p03-w03-p03-047-q0-beth-failure}{Beth and interpolation fail for \(\mathscr L(Q_0)\)}
Let \(E\) be a unary relation symbol, let \(\theta\) be the sentence of
\cref{thm:p03-w03-p03-047-q0-omega}, and let \(\varphi\) be the
\(\mathscr L(Q_0)\)-sentence of vocabulary \(\{<,E\}\)
\[
 \theta\;\land\;\forall x\,\bigl(\forall y\,\neg\,y<x\to Ex\bigr)
 \;\land\;\forall x\,\forall y\,\bigl(S(x,y)\to(Ey\leftrightarrow\neg Ex)\bigr).
\]
\begin{enumerate}
 \item \(E\) is implicitly defined by \(\varphi\) in the sense of
       \cref{def:p03-w03-p03-047-implicit-explicit}: in every model of
       \(\varphi\), \(E\) is the image of the even numbers under the unique
       isomorphism with \((\omega,<)\).
 \item \(E\) is not explicitly definable relative to \(\varphi\) by any
       \(\mathscr L(Q_0)\)-formula \(\psi(c)\) of vocabulary \(\{<,c\}\).
 \item Consequently \(\mathscr L(Q_0)\) lacks the Beth property and, by
       \cref{thm:p03-w03-p03-047-craig-beth}, lacks the interpolation
       property.
\end{enumerate}
\end{hbtheorem}

The proof shows more than part~2 states: over \((\omega,<)\) every
\(\mathscr L(Q_0)\)-formula is equivalent to a first-order formula, and a
first-order formula without parameters cannot separate the even from the odd
numbers.  The implicit definition succeeds because \(\theta\) pins the order
down to \((\omega,<)\), where the recursion along successors has a unique
solution; the explicit definition fails because \(\mathscr L(Q_0)\) adds
nothing over \((\omega,<)\) to what first-order logic says there.  In weak
second-order logic and in \(\omega\)-logic the set of even numbers \emph{is}
explicitly definable over \((\omega,<)\), so those logics need a different
witness.  Mostowski's argument, as Barwise reports it
\cite[Ch.~I, \S1.6, p.~10]{Barwise1985BackgroundAims}, supplies one: a logic
whose formulas are finite objects, codable by numbers, and in which finiteness
is definable can pin down the standard model of arithmetic and then state
the defining clauses of its own satisfaction relation over that model.  The
satisfaction relation is then implicitly defined, and a diagonal argument
shows that it is not explicitly definable.  The next definition isolates
exactly what the argument uses.

\begin{hbdefinition}{def:p03-w03-p03-047-self-satisfaction}{Satisfaction over \(\mathcal N\) defined in the logic}
Let \(\mathscr L\) be a one-sorted logic, \(\tau_<=\Lang_A\cup\{<\}\) as in
\cref{thm:p03-w03-p03-047-q0-not-re}, \(c\) a new constant and \(U\) a new
unary relation symbol.  \(\mathscr L\) \emph{defines its satisfaction over
\(\mathcal N\)} if there are
\begin{enumerate}
 \item[(a)] a set \(\mathrm{Fm}\subseteq\omega\) and a map
       \(e\mapsto\psi_e\) from \(\mathrm{Fm}\) into
       \(\mathscr L[\tau_<\cup\{c\}]\) such that every sentence of
       \(\mathscr L[\tau_<\cup\{c\}]\) has the same models as some
       \(\psi_e\) (\emph{finite syntax coded in \(\omega\)});
 \item[(b)] a function \(g:\omega\times\omega\to\omega\) whose graph is
       defined in \(\mathcal N\) by a first-order formula
       \(\gamma(x,y,z)\) of \(\Lang_A\); and
 \item[(c)] a sentence \(\sigma\in\mathscr L[\tau_<\cup\{U\}]\) such that
       \(U\) is implicitly defined by \(\sigma\)
       (\cref{def:p03-w03-p03-047-implicit-explicit}) and
       \((\mathcal N,<,U^\ast)\models\sigma\) for a set \(U^\ast\subseteq\omega\)
       with \(g(e,n)\in U^\ast\) iff \((\mathcal N,<,n)\models\psi_e\), for all
       \(e\in\mathrm{Fm}\) and \(n\in\omega\).
\end{enumerate}
\end{hbdefinition}

Clause~(c) is where the definability of finiteness enters: uniqueness of
\(U\) needs a sentence of \(\mathscr L\) whose models are copies of
\((\mathcal N,<)\), and in the instances below that sentence comes from the
definition of finiteness (\cref{thm:p03-w03-p03-047-q0-omega} and part~3 of
\cref{thm:p03-w03-p03-047-weak-sol}).

\begin{hbtheorem}{thm:p03-w03-p03-047-mostowski-general}{proved=annex:p03-w03-p03-047-mostowski-general}{Mostowski's theorem: finite syntax and definable finiteness defeat Beth}
Let \(\mathscr L\) be a one-sorted logic
(\cref{def:p03-w03-p03-047-ebbinghaus-logic}) with the atom, Boolean and
particularization properties (\cref{def:p03-w03-p03-047-closure}), and let
\(\mathscr L\) define its satisfaction over \(\mathcal N\) with data
\(\mathrm{Fm},\psi,g,\sigma,U^\ast\).
\begin{enumerate}
 \item \(U\) is implicitly defined by \(\sigma\) but is not explicitly
       definable relative to \(\sigma\).  So \(\mathscr L\) lacks the Beth
       property, and by \cref{thm:p03-w03-p03-047-craig-beth} it lacks the
       interpolation property.
 \item \(\mathscr L(Q_0)\) and one-sorted weak second-order logic
       \(\mathscr L^{w2}\) define their satisfaction over \(\mathcal N\), so
       part~1 applies to both.
\end{enumerate}
\end{hbtheorem}

For \(\mathscr L(Q_0)\) part~2 gives a second proof of the failure in
\cref{thm:p03-w03-p03-047-q0-beth-failure}, with a different witness.
Barwise's statement covers \(\omega\)-logic as well
\cite[Ch.~I, \S1.6, p.~10]{Barwise1985BackgroundAims}.  Read as in
\cref{thm:p03-w03-p03-047-omega-logic-renaming}, \(\omega\)-logic is not a
logic in the sense of \cref{def:p03-w03-p03-047-ebbinghaus-logic}, so
part~1 does not apply to it as stated, and no proof for \(\omega\)-logic is
given here.  The definition above does not cover logics with infinitely long
formulas: in \(\mathscr L_{\omega_1\omega}\) there are uncountably many
sentences, clause~(a) fails, and the Beth property holds
(\cref{thm:p02-w02c-p02-003-beth}).  Barwise draws the moral that a logic
in which finiteness is expressible and which is closed under implicit
definability needs a syntax that is infinitary in some sense
\cite[Ch.~I, \S1.6, p.~10]{Barwise1985BackgroundAims}; the theorem proves
this for the logics meeting the definition above, not for every logic.

\begin{hbtheorem}{thm:p03-w03-p03-047-sol-failures}{proved=annex:p03-w03-p03-047-sol-failures}{Compactness and L\"owenheim--Skolem fail for second-order logic}
Second-order logic \(\mathscr L^2\), with quantification over relations of
every finite arity on the domain, has a sentence of the empty vocabulary true
exactly in the infinite structures and a sentence of the empty vocabulary
true exactly in the uncountable structures.  Hence \(\mathscr L^2\) has
neither the countable compactness property nor the L\"owenheim property.
\end{hbtheorem}

Barwise lists completeness, compactness, interpolation and
L\"owenheim--Skolem as theorems that do not extend to second-order logic
\cite[Ch.~I, \S1.7, p.~11]{Barwise1985BackgroundAims}.  For interpolation this
needs a qualification that Ebbinghaus supplies: in the one-sorted case
\(\mathscr L^2\) has the interpolation property trivially, and it fails only
in the two-sorted case
\cite[Ch.~II, \S2.1, Warning, p.~33]{Ebbinghaus1985ExtendedLogics}.  The
one-sorted claim is proved in \cref{thm:p03-w03-p03-047-sol-interpolation}.
Barwise further notes that monadic second-order logic, which quantifies only
over sets, is weaker than full second-order logic yet supports decidability
results such as Rabin's theorem on the monadic theory of the infinite binary
tree, and that Shelah showed there are only four first-order definable
second-order quantifiers \cite[Ch.~I, \S1.7, p.~12]{Barwise1985BackgroundAims}.

\subsection{Applications}

Barwise separates three kinds of application: outright solutions of open
mathematical questions, independence results (about the limitations of
first-order set theory rather than about extended model theory), and new
concepts that let the right questions be asked
\cite[Ch.~I, \S1.8, pp.~12--13]{Barwise1985BackgroundAims}.  His example of the
third kind is the theorem of Erd\H{o}s, Gillman and Henriksen that two
real-closed fields of order type \(\eta_1\) and cardinality \(\aleph_1\) are
isomorphic.  Without the continuum hypothesis no such fields exist, so this
form of the theorem is vacuous; the reformulation with partial isomorphisms
asserts that any two real-closed fields of order type \(\eta_1\), of any
cardinality, are strongly partially isomorphic
\cite[Ch.~I, \S1.8, p.~13]{Barwise1985BackgroundAims}.  The relevant notion of
potential isomorphism is \cref{def:p03-potential-isomorphism}.

\section{Abstract model theory}
\label{sec:p03-w03-p03-047-barwise-abstract}

\subsection{Lindstr\"om's theorem in Barwise's form}

Barwise states three guiding principles for finding natural logics: build
into the semantics important notions from a domain of mathematics, keep the
semantics constrained to those notions and what they imply, and find a syntax
in which they are naturally expressed \cite[Ch.~I, \S2.1, p.~14]{Barwise1985BackgroundAims}.
He then characterizes first-order logic as the strongest logic with:
(1) the countable compactness property, in the form that a countable set of
sentences without a model has a finite subset without a model; and (2) the
\emph{L\"owenheim property}, that a sentence with an infinite model has a
countable model \cite[Ch.~I, \S2.1, p.~14]{Barwise1985BackgroundAims}.  Here
``countable'' is read as ``of power at most \(\aleph_0\)''.

\begin{hbtheorem}{thm:p03-w03-p03-047-lindstrom-bridge}{proved=annex:p03-w03-p03-047-lindstrom-bridge}{Barwise's hypotheses are those of \cref{thm:p03-lindstrom}}
For a regular abstract logic in the sense of \cref{sec:p03-lindstrom}:
\begin{enumerate}
 \item Barwise's property~(1) is equivalent to hypothesis~1 of
       \cref{thm:p03-lindstrom};
 \item Barwise's property~(2) is equivalent to hypothesis~2 of
       \cref{thm:p03-lindstrom}, that every satisfiable sentence has a model
       of power at most \(\aleph_0\).
\end{enumerate}
Hence a regular logic with (1) and (2) has every sentence equivalent to a
first-order sentence, by \cref{thm:p03-lindstrom}.
\end{hbtheorem}

Both hypotheses are needed: \(\mathscr L(Q_1)\) has~(1) but not~(2), and
\(\mathscr L_{\omega_1\omega}\) has~(2) but not~(1)
\cite[Ch.~I, \S2.1, p.~14]{Barwise1985BackgroundAims}.  The failure of~(2) for
\(\mathscr L(Q_1)\), the failure of~(1) and the validity of~(2) for
\(\mathscr L_{\omega_1\omega}\) are parts~4, 3 and~1 of
\cref{thm:p02-w02-p02-005-bw-extended-ls}; the countable compactness of
\(\mathscr L(Q_1)\) is \cref{thm:p03-keisler-q-one}.  Barwise stresses that Lindstr\"om and
Friedman avoided saying what a logic is by working with classes of
structures closed under conditions, a stop-gap until the notion of logic
itself was analysed \cite[Ch.~I, \S2.1, pp.~14--15]{Barwise1985BackgroundAims};
\cref{def:p03-w03-p03-047-ebbinghaus-logic} is that analysis.

\begin{hbwarning}{def:p03-w03-p03-047-open-1985}{Open questions recorded in 1985 are not theorems}
Barwise records as open, at the time of writing, satisfactory Lindstr\"om-type
characterizations of \(\mathscr L_{\omega_1\omega}\) and of
\(\mathscr L(Q_1)\) \cite[Ch.~I, \S2.2, p.~15]{Barwise1985BackgroundAims}, and
\emph{Feferman's problem}: whether some countably compact logic extending
\(\mathscr L(Q_1)\) has the interpolation property
\cite[Ch.~I, \S2.4, p.~17]{Barwise1985BackgroundAims}.  These are reported as
open questions of the source; nothing in this chapter answers them or treats
them as results.
\end{hbwarning}

\subsection{Interpolation, separation, and \(\Delta\)-closure}

Barwise uses abstract model theory to separate properties that coincide for
first-order logic.  In \(\mathscr L_{\omega_1\omega}\) interpolation holds
(\cref{thm:p02-w02c-p02-003-interpolation}) while Robinson's consistency
property fails; and strong compactness is much
stronger than countable compactness in general, the difference being hidden
for first-order logic by the L\"owenheim--Skolem theorem
\cite[Ch.~I, \S2.3, p.~16]{Barwise1985BackgroundAims}.  The second contrast is
proved for \(\mathscr L(Q_1)\) in \cref{thm:p03-w03-p03-047-q1-not-aleph1}.

\begin{hbdefinition}{def:p03-w03-p03-047-pc-delta}{PC classes and \(\Delta\)-closed logics}
Let \(\mathscr L\) be a logic and \(\tau\) a vocabulary.  A class \(K\) of
\(\tau\)-structures is \emph{PC} (or \(\Sigma^1_1\)) in \(\mathscr L\) if
there are a vocabulary \(\tau'\supseteq\tau\) and a sentence
\(\varphi\in\mathscr L[\tau']\) such that \(\mathfrak M\in K\) iff some
expansion of \(\mathfrak M\) to \(\tau'\) satisfies \(\varphi\).  \(K\) is
\emph{co-PC} if its complement in \(\operatorname{Str}[\tau]\) is PC.  The
logic is \emph{\(\Delta\)-closed} if every class that is both PC and co-PC in
\(\mathscr L\) is the model class of a sentence of \(\mathscr L[\tau]\)
\cite[Ch.~I, \S2.4, pp.~17--18]{Barwise1985BackgroundAims}.  The
\emph{separation form} of interpolation says: whenever two PC classes
\(K_0,K_1\) of \(\tau\)-structures are disjoint, some
\(\theta\in\mathscr L[\tau]\) has \(K_0\subseteq\operatorname{Mod}(\theta)\)
and \(K_1\cap\operatorname{Mod}(\theta)=\varnothing\).
\end{hbdefinition}

\begin{hbtheorem}{thm:p03-w03-p03-047-craig-separation}{proved=annex:p03-w03-p03-047-craig-separation}{Interpolation, separation and \(\Delta\)-closure}
Let \(\mathscr L\) be a one-sorted logic in the sense of
\cref{def:p03-w03-p03-047-ebbinghaus-logic} with the negation property.
\begin{enumerate}
 \item \(\mathscr L\) has the interpolation property of
       \cref{def:p03-w03-p03-047-ebbinghaus-properties} iff it has the
       separation form of \cref{def:p03-w03-p03-047-pc-delta}.
 \item If \(\mathscr L\) has the interpolation property, it is
       \(\Delta\)-closed.
\end{enumerate}
\end{hbtheorem}

The converse of the second clause fails in general: Barwise states that
\(\Delta\)-closure is strictly weaker than interpolation, that every logic has
a least \(\Delta\)-closed extension \(\Delta(\mathscr L)\), and that
\(\Delta(\mathscr L(Q_0))\), the common \(\Delta\)-closure of
\(\mathscr L(Q_0)\), weak second-order logic and \(\omega\)-logic, is the
hyperarithmetic fragment of \(\mathscr L_{\omega_1\omega}\), a result of
Makowsky and of Barwise \cite[Ch.~I, \S2.4, p.~18]{Barwise1985BackgroundAims}.
Neither \(\mathscr L(Q_0)\) nor \(\mathscr L(Q_1)\) has the interpolation
property; for \(\mathscr L(Q_0)\) the reason is Mostowski's (an explicit
instance is \cref{thm:p03-w03-p03-047-q0-beth-failure}), while for
\(\mathscr L(Q_1)\) the counterexamples suggested missing notions near
countability rather than an essential obstacle
\cite[Ch.~I, \S2.4, pp.~16--17]{Barwise1985BackgroundAims}.  Barwise also
reports that a logic whose sentences have finitely many, or reasonably few,
symbols has Robinson's consistency property just when it is compact and has
the interpolation property \cite[Ch.~I, \S2.4, p.~16]{Barwise1985BackgroundAims}.

\subsection{Pinning down ordinals and Hanf numbers}

Barwise calls a logic \emph{bounded by} \(\alpha\) when \(\alpha\) exceeds
every ordinal that can be characterized in it, and \emph{unbounded} when
well-ordering is definable \cite[Ch.~I, \S2.5, p.~18]{Barwise1985BackgroundAims}.
Ebbinghaus makes this exact with the notions of a set of sentences pinning
down an ordinal \(\alpha\) via \(<\), and of the well-ordering numbers
\(w_\kappa(\mathscr L)\) and \(w(\mathscr L)=w_1(\mathscr L)\), fixed in
\cref{def:p03-w03-p03-041-pinning}.

\begin{hbdefinition}{def:p03-w03-p03-047-bounded}{Bounded logic}
A logic \(\mathscr L\) is \emph{bounded} if no single \(\mathscr L\)-sentence
pins down arbitrarily large ordinals, in the sense of
\cref{def:p03-w03-p03-041-pinning}
\cite[Ch.~II, Definition~5.2.1, p.~61]{Ebbinghaus1985ExtendedLogics}.
Barwise's phrase ``bounded by \(\alpha\)'' says that every ordinal pinned
down by a sentence is below \(\alpha\).
\end{hbdefinition}

\begin{hbtheorem}{thm:p03-w03-p03-047-bounded-omega}{proved=annex:p03-w03-p03-047-bounded-omega}{Countably compact regular logics are bounded by \(\omega\)}
Let \(\mathscr L\) be a one-sorted logic with the atom, Boolean and
particularization properties of \cref{def:p03-w03-p03-047-closure}, and with
the countable compactness property.  Then for each \(\mathscr L\)-sentence
there is a natural number above every ordinal it pins down; every finite
ordinal \(n\ge 2\) is pinned down by a first-order sentence; hence
\(\mathscr L\) is bounded and \(w(\mathscr L)=\omega\).  In particular
\(w(\mathscr L_{\omega\omega})=\omega\), and \(w(\mathscr L(Q_1))=\omega\)
given the countable compactness of \(\mathscr L(Q_1)\)
(\cref{thm:p03-keisler-q-one}).
\end{hbtheorem}

Barwise reports in addition that \(\mathscr L(Q_0)\) is bounded by
\(\omega_1^{c}\), the least non-recursive ordinal, that
\(\mathscr L_{\omega_1\omega}\) is bounded by \(\omega_1\), that the
second-order-like logics \(\mathscr L_{\omega_1\omega_1}\) and the logic of the
game quantifier are not bounded, and that first-order logic is
characterized by the L\"owenheim--Skolem theorem together with boundedness by
\(\omega\) \cite[Ch.~I, \S2.5, p.~18]{Barwise1985BackgroundAims}.

\begin{hbdefinition}{def:p03-w03-p03-047-hanf}{Hanf number; the beth function}
If the sentences of a logic \(\mathscr L\) form a set, the \emph{Hanf
number} \(h(\mathscr L)\) is the least cardinal \(\kappa\) such that every
sentence of \(\mathscr L\) with a model of power at least \(\kappa\) has
models of arbitrarily large power.  For a countable theory \(T\),
\(h(T,\mathscr L)\) is the least \(\kappa\) such that for every
\(\mathscr L\)-sentence \(\varphi\), if \(T\cup\{\varphi\}\) has a model of
power \(\kappa\) then it has arbitrarily large models.  The beth numbers are
\(\beth_0=\aleph_0\), \(\beth_{\alpha+1}=2^{\beth_\alpha}\), and
\(\beth_\lambda=\sup_{\alpha<\lambda}\beth_\alpha\) for limit \(\lambda\)
\cite[Ch.~I, \S2.6, pp.~19--20]{Barwise1985BackgroundAims}.
\end{hbdefinition}

\begin{hbtheorem}{thm:p03-w03-p03-047-hanf-exists}{proved=annex:p03-w03-p03-047-hanf-exists}{Hanf's observation}
If the sentences of \(\mathscr L\) form a set, then \(h(\mathscr L)\) exists.
For first-order logic in a fixed countable vocabulary,
\(h(\mathscr L_{\omega\omega})=\aleph_0\).
\end{hbtheorem}

The proof is non-constructive, as Barwise emphasizes: it bounds
\(h(\mathscr L)\) by a supremum over the set of sentences without computing
it \cite[Ch.~I, \S2.6, p.~19]{Barwise1985BackgroundAims}.  For
\(\mathscr L(Q_0)\), \(\mathscr L_{\omega_1\omega}\) and
\(\mathscr L_{\kappa\omega}\) Barwise states that \(h(\mathscr L)=\beth_\lambda\)
where \(\lambda\) is the least ordinal not pinned down by the logic, and that
for unbounded logics little can be said; Shelah's structural explanation,
in which a theory is bounded exactly when its models decompose as products
of countable models and exactly when its Hanf number is easily computed, is
deferred by Barwise to a later chapter of the volume
\cite[Ch.~I, \S2.6, pp.~19--20]{Barwise1985BackgroundAims}.  First-order logic
shows that the formula is not meant for every logic:
\(h(\mathscr L_{\omega\omega})=\aleph_0\) by \cref{thm:p03-w03-p03-047-hanf-exists}.

\subsection{Large cardinals, other structures, unnatural logics}

Barwise records that weakly and strongly compact cardinals \(\kappa\) are
defined by analogues of countable and full compactness for
\(\mathscr L_{\kappa\kappa}\); that measurable cardinals arise with
Robinson's consistency property; and that Vop\v{e}nka's principle holds exactly when each finitely generated logic
possesses a \emph{strong compactness cardinal}, a \(\kappa\) such that every theory
without a model has a subset of power less than \(\kappa\) without a model
\cite[Ch.~I, \S2.7, pp.~20--21]{Barwise1985BackgroundAims}.  He presents these
as results about very abstract logics rather than about logics arising from
mathematical concepts.  Lindstr\"om's dilemma (give up compactness or
L\"owenheim--Skolem) is escaped by changing the structures, as in
topological, measure and probability logics, whose settings violate
assumptions usually left implicit \cite[Ch.~I, \S2.8, p.~21]{Barwise1985BackgroundAims}.

\begin{hbwarning}{def:p03-w03-p03-047-unnatural}{Counterexample logics}
Some logics exist only to show that one property of logics does not imply
another, and some logics that looked natural, \(\mathscr L(Q_0)\) among them
in Barwise's judgement, turned out not to be
\cite[Ch.~I, \S2.9, pp.~21--22]{Barwise1985BackgroundAims}.  Such logics test the
relations between properties; they are not evidence about the logic of a
mathematical concept.
\end{hbwarning}

\subsection{Barwise's conclusion}

Barwise's closing section states limitations rather than results
\cite[Ch.~I, \S3, pp.~22--23]{Barwise1985BackgroundAims}.  Work that belongs
under the heading of model-theoretic logics but sits badly in the general
framework of abstract model theory is, by his account, not usually counted as
extended model theory, and the volume does not attempt to include it
\cite[Ch.~I, \S3, p.~22]{Barwise1985BackgroundAims}.  The largest such omission
is the semantics of computer languages, which would need a separate volume;
that work exposes a weakness of traditional model theory, its concern with
objects and static relations among them rather than with activity, and the
same weakness troubles the traditional attempts to apply model theory to
human languages, a subject the volume also leaves aside
\cite[Ch.~I, \S3, p.~22]{Barwise1985BackgroundAims}.  Barwise expects that
in the long run a unified view of logic and semantics will need one
framework covering both the model theory of mathematical logics and the
semantics of computer languages, and he records that no such framework was
in sight when he wrote \cite[Ch.~I, \S3, p.~22]{Barwise1985BackgroundAims}.

Three further limitations are named \cite[Ch.~I, \S3, p.~22]{Barwise1985BackgroundAims}.
First, traditional model theory is about the truth and satisfaction of
sentences, so it has no place for the use of language to bring about a
change.  Barwise credits this criticism to Austin and the other writers on
natural language in the speech-act tradition, and he locates the same
phenomenon in computer languages, where the power of a statement to change
the state of the system (its \emph{side effects}) is what most separates
their semantics from that of a logic in the sense of
\cref{def:p03-w03-p03-047-barwise-logic}.  Second, inference, whether
carried out by a person or a machine, is an activity of extracting
information, while the tradition reduces it to objects such as proofs and
strings of symbols; Barwise proposes elsewhere to treat traditional
inference as one kind of information-preserving activity and leaves that
development to other work.  Third, even inference in the traditional sense
has made little progress in extended model theory, which has few genuine
completeness theorems (\cref{thm:p03-w03-p03-047-abstract-completeness}) and
fewer extensions of proof theory.

His positive programme is on p.~23 \cite[Ch.~I, \S3, p.~23]{Barwise1985BackgroundAims}.
Mathematicians lose patience with logic, he says, because many of the notions
they use are not first-order while first-order logic has been presented to
them as logic itself.  The study of model-theoretic logics is to change this
in three ways: by finding the logic of the concepts mathematicians actually
use (finiteness, countability, well-ordering, and the others of
\cref{sec:p03-w03-p03-047-abstract-logics}), by finding applications, and by
isolating new concepts that enrich both mathematics and logic.  Barwise
states that he does not know how much of the work in the volume will keep a
permanent place in mathematics, the subject being young; what he treats as
settled, whatever happens to the particular results, is that the
identification of logic with first-order logic
(\cref{def:p03-w03-p03-047-first-order-thesis}) cannot be restored.  These
are programmatic claims about the scope of the subject; they carry no
mathematical hypothesis to check, and the uncertainty about which
particulars will last is Barwise's own.

\section{The general framework}
\label{sec:p03-w03-p03-047-ebbinghaus-framework}

Ebbinghaus fixes the notion of logic used in the rest of the volume.  He
states it for many-sorted vocabularies; examples and proofs are one-sorted
unless stated otherwise \cite[Ch.~II, introduction, p.~25]{Ebbinghaus1985ExtendedLogics}.

\begin{hbconvention}{def:p03-w03-p03-047-vocabularies}{Vocabularies, structures, reducts, renamings}
A many-sorted \emph{vocabulary} \(\tau\) is a nonempty set of sort symbols,
finitary relation symbols, finitary function symbols and constants; each
constant, each function symbol and each argument place carries a sort symbol
of \(\tau\).  A \emph{\(\tau\)-structure} \(\mathfrak A\) has a nonempty
domain \(A_s\) for each sort \(s\) and interprets the remaining symbols in
the usual way; \(\operatorname{Str}[\tau]\) is the class of
\(\tau\)-structures and \(\tau_{\mathfrak A}\) the vocabulary of
\(\mathfrak A\).  Equations are formed only between terms of the same sort.
For \(\sigma\subseteq\tau\), the \emph{reduct} \(\mathfrak A\restriction\sigma\)
forgets the domains of sorts outside \(\sigma\) and the interpretations of
symbols outside \(\sigma\).  In the one-sorted case, \(C\subseteq A\) is
\(\tau\)-closed if it is nonempty, contains every constant and is closed
under every function; \(\mathfrak A\,|\,C\) is then the substructure on
\(C\), and \((\mathfrak A\restriction\sigma)\,|\,P^{\mathfrak A}\) is a
\emph{relativized reduct}.  A \emph{renaming} \(\rho:\tau\to\sigma\) is a
bijection sending sorts to sorts, relation and function symbols to symbols
of the same kind and arity, and constants to constants, compatibly with the
attached sorts; \(\mathfrak A^\rho\) is the \(\sigma\)-structure with
\(B_{\rho(s)}=A_s\) and \(\rho(S)^{\mathfrak A^\rho}=S^{\mathfrak A}\)
\cite[Ch.~II, \S1.1, pp.~26--27]{Ebbinghaus1985ExtendedLogics}.
\end{hbconvention}

\begin{hbexample}{def:p03-w03-p03-047-sorted-vocabulary}{A three-sorted vocabulary}
Let \(R\) be binary with argument places of sorts \(s_2,s_1\) in that order,
\(f\) unary with argument of sort \(s_1\) and value of sort \(s_2\), and \(c\)
a constant of sort \(s_1\).  Then \(\tau=\{s_1,s_2,s_3,R,f,c\}\) is a
vocabulary.  The term \(f(c)\) has sort \(s_2\), but \(f(f(c))\) is not a
term, since \(f\) needs an argument of sort \(s_1\).  Atomic sentences are
\(Rt_0t_1\) with \(t_0\) of sort \(s_2\) and \(t_1\) of sort \(s_1\), and
equations between terms of one sort.  A \(\tau\)-structure has three nonempty
domains, a relation \(R^{\mathfrak A}\subseteq A_{s_2}\times A_{s_1}\), a map
\(f^{\mathfrak A}:A_{s_1}\to A_{s_2}\) and an element
\(c^{\mathfrak A}\in A_{s_1}\); its reduct to \(\{s_1,s_2,R\}\) forgets
\(A_{s_3}\), \(f^{\mathfrak A}\) and \(c^{\mathfrak A}\)
\cite[Ch.~II, \S1.1, pp.~26--27]{Ebbinghaus1985ExtendedLogics}.  The sort
\(s_3\) carries no symbol, yet its domain is part of every \(\tau\)-structure
and counts towards the power defined in
\cref{def:p03-w03-p03-047-ebbinghaus-properties}.
\end{hbexample}

\begin{hbdefinition}{def:p03-w03-p03-047-ebbinghaus-logic}{Logic (Ebbinghaus, Definition~1.1.1)}
A \emph{logic} is a pair \((\mathscr L,\models_{\mathscr L})\) where
\(\mathscr L\) assigns to each vocabulary \(\tau\) a class
\(\mathscr L[\tau]\) (the \(\mathscr L\)-sentences of vocabulary \(\tau\))
and \(\models_{\mathscr L}\) is a relation between structures and
\(\mathscr L\)-sentences, such that:
\begin{enumerate}
 \item[(i)] if \(\tau\subseteq\sigma\) then
       \(\mathscr L[\tau]\subseteq\mathscr L[\sigma]\);
 \item[(ii)] if \(\mathfrak A\models_{\mathscr L}\varphi\) then
       \(\varphi\in\mathscr L[\tau_{\mathfrak A}]\);
 \item[(iii)] \emph{isomorphism property}: if
       \(\mathfrak A\models_{\mathscr L}\varphi\) and
       \(\mathfrak B\cong\mathfrak A\) then
       \(\mathfrak B\models_{\mathscr L}\varphi\);
 \item[(iv)] \emph{reduct property}: if \(\varphi\in\mathscr L[\tau]\) and
       \(\tau\subseteq\tau_{\mathfrak A}\) then
       \(\mathfrak A\models_{\mathscr L}\varphi\) iff
       \(\mathfrak A\restriction\tau\models_{\mathscr L}\varphi\);
 \item[(v)] \emph{renaming property}: for each renaming
       \(\rho:\tau\to\sigma\) and \(\varphi\in\mathscr L[\tau]\) there is
       \(\varphi^\rho\in\mathscr L[\sigma]\) with
       \(\mathfrak A\models_{\mathscr L}\varphi\) iff
       \(\mathfrak A^\rho\models_{\mathscr L}\varphi^\rho\) for all
       \(\tau\)-structures \(\mathfrak A\)
\end{enumerate}
\cite[Ch.~II, Definition~1.1.1, pp.~27--28]{Ebbinghaus1985ExtendedLogics}.
Write \(\operatorname{Mod}^\tau_{\mathscr L}(\varphi)=\{\mathfrak A\in
\operatorname{Str}[\tau]:\mathfrak A\models_{\mathscr L}\varphi\}\); for
\(\Phi\cup\{\varphi\}\subseteq\mathscr L[\tau]\),
\(\Phi\models_{\mathscr L}\varphi\) iff every \(\tau\)-structure satisfying
all of \(\Phi\) satisfies \(\varphi\); \(\mathfrak A\equiv_{\mathscr L}\mathfrak B\)
iff the two \(\tau\)-structures satisfy the same sentences of
\(\mathscr L[\tau]\); and \(\operatorname{Th}_{\mathscr L}(\mathfrak A)\) is
the set of sentences of \(\mathscr L[\tau_{\mathfrak A}]\) true in
\(\mathfrak A\) \cite[Ch.~II, \S1.1, p.~28]{Ebbinghaus1985ExtendedLogics}.
No rule of inference is part of the definition; Ebbinghaus notes that
``model-theoretic language'' would be the more accurate name.
\end{hbdefinition}

Ebbinghaus allows two variants.  In the first, a logic carries variables of
each sort and two assignments \(\operatorname{Form}_{\mathscr L}\) and
\(\operatorname{Sent}_{\mathscr L}\); since free variables behave like
constants the difference is minor \cite[Ch.~II, \S1.1.2, pp.~28--29]{Ebbinghaus1985ExtendedLogics}.
First-order logic \(\mathscr L_{\omega\omega}\), second-order logic
\(\mathscr L^2\), weak second-order logic \(\mathscr L^{w2}\),
\(\mathscr L_{\omega_1\omega}\), \(\mathscr L_{\infty\omega}\) and
\(\mathscr L_{\omega\omega}(Q_\alpha)\) are logics in both senses
\cite[Ch.~II, \S1.1.2, p.~29]{Ebbinghaus1985ExtendedLogics}.  In the second
variant a logic also carries a class \(\operatorname{Str}_{\mathscr L}[\tau]\)
of permitted \(\tau\)-structures for each \(\tau\), the properties
(i)--(v) being modified accordingly; this is needed for \(\omega\)-logic,
which admits only structures \(\mathfrak A\) with
\(\{s,<\}\subseteq\tau_{\mathfrak A}\) and
\(\mathfrak A\restriction\{s,<\}\cong(\omega,<)\)
\cite[Ch.~II, \S1.1.3, p.~29]{Ebbinghaus1985ExtendedLogics}.

\begin{hbtheorem}{thm:p03-w03-p03-047-omega-logic-renaming}{proved=annex:p03-w03-p03-047-omega-logic-renaming}{\(\omega\)-logic violates the renaming property}
Read \(\omega\)-logic within \cref{def:p03-w03-p03-047-ebbinghaus-logic} as
follows, with the single sort \(s\): \(\mathscr L[\tau]\) is the set of
first-order sentences of \(\tau\) when \(\{s,<\}\subseteq\tau\) and is empty
otherwise, and \(\mathfrak A\models\varphi\) iff
\(\varphi\in\mathscr L[\tau_{\mathfrak A}]\),
\(\mathfrak A\restriction\{s,<\}\cong(\omega,<)\), and \(\mathfrak A\)
satisfies \(\varphi\) in the first-order sense.  Then:
\begin{enumerate}
 \item properties (i)--(iv) of \cref{def:p03-w03-p03-047-ebbinghaus-logic}
       hold and the renaming property (v) fails;
 \item the failure does not depend on the empty sentence classes: if
       \(\rho:\{s,<\}\to\{s,\prec\}\) renames \(<\) to a fresh binary
       \(\prec\), no first-order sentence \(\psi\) of \(\{s,\prec\}\),
       evaluated by ordinary first-order satisfaction, has
       \(\mathfrak A\models\forall x\,x=x\) iff \(\mathfrak A^\rho\models\psi\)
       for all \(\{s,<\}\)-structures \(\mathfrak A\).
\end{enumerate}
\end{hbtheorem}

\subsection{Closure properties and regular logics}

\begin{hbdefinition}{def:p03-w03-p03-047-closure}{Basic closure properties (Ebbinghaus, Definition~1.2.1)}
A logic \(\mathscr L\) has:
\begin{enumerate}
 \item[(i)] the \emph{atom property} if for every \(\tau\) and every atomic
       \(\varphi\in\mathscr L_{\omega\omega}[\tau]\) some
       \(\psi\in\mathscr L[\tau]\) has
       \(\operatorname{Mod}^\tau_{\mathscr L}(\psi)=\operatorname{Mod}^\tau_{\mathscr L_{\omega\omega}}(\varphi)\);
 \item[(ii)] the \emph{negation property} if for every
       \(\varphi\in\mathscr L[\tau]\) some \(\psi\in\mathscr L[\tau]\) has
       \(\operatorname{Mod}^\tau_{\mathscr L}(\psi)=\operatorname{Str}[\tau]\setminus\operatorname{Mod}^\tau_{\mathscr L}(\varphi)\);
 \item[(iii)] the \emph{conjunction property} if for all
       \(\varphi_0,\varphi_1\in\mathscr L[\tau]\) some
       \(\psi\in\mathscr L[\tau]\) has
       \(\operatorname{Mod}^\tau_{\mathscr L}(\psi)=\operatorname{Mod}^\tau_{\mathscr L}(\varphi_0)\cap\operatorname{Mod}^\tau_{\mathscr L}(\varphi_1)\);
 \item[(iv)] the \emph{particularization property} if for every constant
       \(c\in\tau\) of sort \(s\) and every \(\varphi\in\mathscr L[\tau]\)
       some \(\psi\in\mathscr L[\tau\setminus\{c\}]\) has, for all
       \((\tau\setminus\{c\})\)-structures \(\mathfrak A\),
       \(\mathfrak A\models\psi\) iff \((\mathfrak A,a)\models\varphi\) for
       some \(a\in A_s\)
\end{enumerate}
\cite[Ch.~II, Definition~1.2.1, pp.~29--30]{Ebbinghaus1985ExtendedLogics}.
The \emph{Boolean property} is (ii) with (iii); the chosen sentences are
written \(\neg\varphi\), \(\varphi_0\land\varphi_1\) and \(\exists c\,\varphi\).
\end{hbdefinition}

\begin{hbdefinition}{def:p03-w03-p03-047-relativization}{Relativization and substitution (Ebbinghaus, Definitions~1.2.2 and~1.2.3)}
Let \(\mathscr L\) be one-sorted.  It has the \emph{relativization property}
if for \(c\notin\tau\cup\sigma\), \(\chi\in\mathscr L[\sigma\cup\{c\}]\) and
\(\varphi\in\mathscr L[\tau]\) there is \(\psi\in\mathscr L[\tau\cup\sigma]\)
such that for every \((\tau\cup\sigma)\)-structure \(\mathfrak B\) in which
\(\chi^{\mathfrak B}=\{b\in B:(\mathfrak B,b)\models\chi\}\) is
\(\tau\)-closed, \(\mathfrak B\models\psi\) iff
\((\mathfrak B\restriction\tau)\,|\,\chi^{\mathfrak B}\models\varphi\); such
a \(\psi\) is written \(\varphi^{\{c\mid\chi(c)\}}\), or \(\varphi^P\) when
\(\chi=Pc\) \cite[Ch.~II, Definition~1.2.2, p.~30]{Ebbinghaus1985ExtendedLogics}.
It has the \emph{substitution property} if the following holds.
Let \(\tau\subseteq\tau'\), let \(\varphi\in\mathscr L[\tau']\),
and, for each symbol \(S\in\tau'\setminus\tau\) other than the
retained sort, choose a formula \(\psi_S(\mathbf c_S)\) over \(\tau\)
with fresh parameter constants: \(n\) for an \(n\)-ary relation,
\(n+1\) for an \(n\)-ary function, and one for a constant.
There is \(\varphi^\ast\in\mathscr L[\tau]\) with the following
satisfaction property. In every \(\tau\)-structure \(\mathfrak A\)
where all prescribed function predicates are total single-valued graphs
and all prescribed constant predicates are singletons, interpret the
new symbols by those graphs and by the relations defined by the
\(\psi_S\). For the unique resulting expansion \(\mathfrak A'\),
\[
 \mathfrak A\models\varphi^\ast
 \quad\Longleftrightarrow\quad
 \mathfrak A'\models\varphi.
\]
Such a structure is \emph{admissible for the substitution}. No condition
is imposed on \(\varphi^\ast\) at inadmissible structures. This is
simultaneous substitution on admissible interpretations; the family
of replaced symbols need not be finite.

For a logic with Boolean operations and particularization, finitely
many function and constant replacements can be guarded. Their graph
conditions are
\[
 \forall\mathbf x\,\exists!y\,\psi_f(\mathbf x,y),
 \qquad \exists!y\,\psi_c(y).
\]
Quantification and Boolean operations put these conditions in the logic.
Conjoin them with \(\varphi^\ast\). The resulting sentence is false
at every inadmissible structure. For one replacement this recovers
the existential graph descriptions of
\cite[Ch.~II, Definitions~1.2.2--1.2.3, pp.~30--31]{Ebbinghaus1985ExtendedLogics}.
With infinitely many replacements, admissibility is a condition on
interpretations, not a hidden infinitary conjunction in one sentence.
\Cref{thm:p03-finite-occurrence-graph-obstruction} proves why this
distinction is necessary even in first-order logic.
\end{hbdefinition}

A logic with the properties of
\cref{def:p03-w03-p03-047-closure,def:p03-w03-p03-047-relativization} is
\emph{regular}.  Every first-order sentence is matched, inside a regular logic, by
an equivalent sentence over that vocabulary (one with exactly the same models), and it is customary to
identify the two \cite[Ch.~II, \S1.2, p.~31]{Ebbinghaus1985ExtendedLogics}.
The relational package of \cref{sec:p03-lindstrom} is specified
separately: it requires finite occurrence vocabularies and the operations
used in that characterization, rather than all the substitutions above.  A logic has
the \emph{occurrence property} if every sentence \(\varphi\) lies in a
smallest \(\mathscr L[\tau_\varphi]\); for questions of effectiveness one
assumes \(\tau_\varphi\) finite \cite[Ch.~II, \S1.2, p.~31]{Ebbinghaus1985ExtendedLogics}.

\begin{hbtheorem}{thm:p03-w03-p03-047-function-elimination}{proved=annex:p03-w03-p03-047-function-elimination}{Graph elimination}
Let \(\mathscr L\) be a one-sorted regular logic. Replace each function
symbol \(f\) of arity \(r\) in a vocabulary \(\tau\) by a new
relation \(R_f\) of arity \(r+1\), and each constant \(c\) by a new unary relation \(R_c\),
retaining the relation symbols and the sort. Call the resulting
vocabulary \(\sigma\). For a \(\tau\)-structure \(\mathfrak A\),
write \(\mathfrak A^{\mathrm{gr}}\) for its \(\sigma\)-structure of graphs.
Let \(\Gamma_\tau\) consist of
\[
 \begin{aligned}
 &\forall\bar x\,\exists!y\,R_f(\bar x,y)
      &&(f\text{ a function symbol of }\tau),\\
 &\exists!y\,R_c(y)
      &&(c\text{ a constant of }\tau).
 \end{aligned}
\]
For \(\varphi\in\mathscr L[\tau]\) there is
\(\varphi^{\mathrm{gr}}\in\mathscr L[\sigma]\) such that:
\begin{enumerate}
 \item On every \(\sigma\)-structure satisfying \(\Gamma_\tau\),
       \(\varphi^{\mathrm{gr}}\) holds exactly when \(\varphi\)
       holds in the uniquely reconstructed \(\tau\)-structure.
 \item Over all \(\sigma\)-structures, the models of the theory
       \(\Gamma_\tau\cup\{\varphi^{\mathrm{gr}}\}\) are exactly
       the graph structures of models of \(\varphi\).
 \item If \(\tau\) has only finitely many function and constant symbols,
       this class is defined by the single sentence
       \((\bigwedge\Gamma_\tau)\land\varphi^{\mathrm{gr}}\).
\end{enumerate}
There is no restriction on the number of relation symbols.
The finite graph replacement is the function-elimination construction of
\cite[Ch.~II, \S1.2, p.~31]{Ebbinghaus1985ExtendedLogics}; the displayed
theory states its unrestricted-vocabulary form explicitly.
\end{hbtheorem}

\begin{hbtheorem}{thm:p03-finite-occurrence-graph-obstruction}{proved=annex:p03.graph.1}{Graph admissibility}
Let \(\mathscr L\) have the reduct property and finite occurrence:
each sentence belongs to the logic of some finite sub-vocabulary.
Let \(\sigma\) consist of binary relations \(R_n\) for \(n<\omega\).
No \(\mathscr L[\sigma]\)-sentence defines the class in which every
\(R_n\) is the graph of a total unary function.
Thus the unrestricted graph theory of
\cref{thm:p03-w03-p03-047-function-elimination} need not be equivalent
to one first-order sentence.
\end{hbtheorem}

The distinction is between a single defining sentence and a theory of
graph conditions, not between finite and infinite structures. Arbitrary
simultaneous substitutions remain available on admissible interpretations.
In \cref{thm:p03-general-fibres}, every substituted graph is the graph
of an explicit total term. All auxiliary structures are therefore
admissible, even for an uncountable vocabulary; no additional graph
conditions enlarge the countable theory used there.

\begin{hbtheorem}{thm:p03-w03-p03-047-substitution-renaming}{proved=annex:p03-w03-p03-047-substitution-renaming}{Substitution yields renaming in the one-sorted case}
A one-sorted logic with the atom property and the substitution property has
the renaming property \cite[Ch.~II, \S1.2, p.~31]{Ebbinghaus1985ExtendedLogics}.
\end{hbtheorem}

Ebbinghaus adds that the many-sorted version requires a careful treatment of
sort symbols; the proof in the annex uses a single sort throughout, and the
many-sorted case is not claimed here.

\subsection{Properties carried over from first-order logic}

\begin{hbdefinition}{def:p03-w03-p03-047-ebbinghaus-properties}{Compactness, effectiveness, L\"owenheim--Skolem, interpolation, Beth (Ebbinghaus, Definition~1.2.4)}
Let \(\mathscr L\) be a logic.
\begin{enumerate}
 \item[(i)] For an infinite cardinal \(\kappa\), \(\mathscr L\) is
       \emph{\(\kappa\)-compact} if for all \(\tau\) and all
       \(\Phi\subseteq\mathscr L[\tau]\) of power at most \(\kappa\),
       \(\Phi\) is satisfiable whenever all its finite subsets are.
 \item[(ii)] \(\mathscr L\) is \emph{compact} if it is \(\kappa\)-compact
       for every infinite \(\kappa\).
 \item[(iii)] \(\mathscr L\) is \emph{effective} if for all
       \(\tau\subseteq\mathrm{HF}\) (the hereditarily finite sets)
       \(\mathscr L[\tau]=\bigcup\{\mathscr L[\tau_0]:\tau_0\subseteq\tau,\ \tau_0\text{ finite}\}\),
       and for every \(\tau_0\in\mathrm{HF}\), \(\mathscr L[\tau_0]\) is a
       recursive subset of \(\mathrm{HF}\).
 \item[(iv)] \(\mathscr L\) is \emph{effectively regular} if it is regular
       and effective and every regularity property is effective; for
       instance, for each \(\tau_0\in\mathrm{HF}\) there is a recursive
       \(\neg:\mathscr L[\tau_0]\to\mathscr L[\tau_0]\) with
       \(\neg(\varphi)\) a negation of \(\varphi\).
 \item[(v)] \(\mathscr L\) is \emph{recursively enumerable for validity} if
       it is effective and for every \(\tau_0\in\mathrm{HF}\) the set of
       valid sentences of \(\mathscr L[\tau_0]\) is recursively enumerable.
 \item[(vi)] \(\mathscr L\) is \emph{recursively enumerable for
       consequence} if it is effective and for every \(\tau_0\in\mathrm{HF}\)
       and every recursively enumerable \(\Phi\subseteq\mathscr L[\tau_0]\),
       \(\{\varphi\in\mathscr L[\tau_0]:\Phi\models\varphi\}\) is
       recursively enumerable.
 \item[(vii)] \(\mathscr L\) has the \emph{L\"owenheim--Skolem property}
       (\emph{down to \(\kappa\)}) if each satisfiable sentence has a model
       of power at most \(\aleph_0\) (at most \(\kappa\)); the power of a
       structure is \(|A|\) in the one-sorted case and
       \(\sum_{s\in\tau}|A_s|\) in the many-sorted case.
 \item[(viii)] \(\mathscr L\) satisfies \emph{Craig interpolation} (the
       \emph{interpolation property}) if whenever \(\varphi_i\in\mathscr L[\tau_i]\)
       (\(i=0,1\)) and \(\varphi_0\models\varphi_1\), some
       \(\psi\in\mathscr L[\tau_0\cap\tau_1]\) has
       \(\varphi_0\models\psi\) and \(\psi\models\varphi_1\), provided in the
       many-sorted case that \(\tau_0\cap\tau_1\) contains a sort symbol.
 \item[(ix)] \(\mathscr L\) has the \emph{Beth property} if for all
       \(\tau\), all symbols \(S\in\tau\) other than sort symbols and all
       \(\varphi\in\mathscr L[\tau]\), if \(S\) is implicitly defined by
       \(\varphi\) then \(S\) is explicitly definable relative to \(\varphi\)
\end{enumerate}
\cite[Ch.~II, Definition~1.2.4, pp.~31--32]{Ebbinghaus1985ExtendedLogics}.
\end{hbdefinition}

\begin{hbdefinition}{def:p03-w03-p03-047-implicit-explicit}{Implicit and explicit definability (Ebbinghaus, Definition~1.2.5)}
Let \(R\in\tau\) be unary of sort \(s\) and \(\varphi\in\mathscr L[\tau]\).
(i) \(R\) is \emph{implicitly defined} by \(\varphi\) if no
\((\tau\setminus\{R\})\)-structure admits two distinct expansions to
\(\tau\)-structures that satisfy \(\varphi\).  (ii) \(R\) is \emph{explicitly
definable relative to} \(\varphi\) if, for a new constant \(c\) of sort
\(s\), some \(\psi(c)\in\mathscr L[(\tau\setminus\{R\})\cup\{c\}]\) satisfies
\(R^{\mathfrak A}=\{a\in A_s:(\mathfrak A,a)\models\psi(c)\}\) for every
\(\tau\)-structure \(\mathfrak A\models\varphi\); intuitively,
\(\varphi\models\forall c(Rc\leftrightarrow\psi(c))\)
\cite[Ch.~II, Definition~1.2.5, p.~32]{Ebbinghaus1985ExtendedLogics}.  For
an \(n\)-ary \(R\) the same clauses are used with an \(n\)-tuple
\(\mathbf c\) of new constants.
\end{hbdefinition}

The source gives these notions ``say for unary \(R\)''; it does not define
explicit definability for function symbols or constants, and none is
supplied here.  Accordingly the next theorem is stated for relation symbols.

\begin{hbtheorem}{thm:p03-w03-p03-047-craig-beth}{proved=annex:p03-w03-p03-047-craig-beth}{Interpolation implies the Beth property}
Let \(\mathscr L\) be a one-sorted logic with the atom, negation and
conjunction properties and the interpolation property.  If a relation symbol
\(R\in\tau\) is implicitly defined by \(\varphi\in\mathscr L[\tau]\), then
\(R\) is explicitly definable relative to \(\varphi\).  In particular every
regular logic with the interpolation property has the Beth property for
relation symbols \cite[Ch.~II, \S1.2, p.~32]{Ebbinghaus1985ExtendedLogics}.
\end{hbtheorem}

Ebbinghaus calls this the first theorem of abstract model theory and poses
the converse question, under which conditions the Beth property yields
interpolation, which the volume takes up in a later chapter
\cite[Ch.~II, \S1.2, p.~32]{Ebbinghaus1985ExtendedLogics}.  The proof uses
only properties (i)--(v) of \cref{def:p03-w03-p03-047-ebbinghaus-logic} and
the closure properties named, which is Ebbinghaus's point: the usual
first-order argument is already an argument about abstract logics.

\begin{hbtheorem}{thm:p03-w03-p03-047-q1-not-aleph1}{proved=annex:p03-w03-p03-047-q1-not-aleph1}{\(\mathscr L(Q_1)\) is not \(\aleph_1\)-compact}
In the vocabulary \(\{c_\alpha:\alpha<\aleph_1\}\) of distinct constants,
the set
\[
 \{\neg Q_1x\,(x=x)\}\cup\{\neg c_\alpha=c_\beta:\alpha<\beta<\aleph_1\}
\]
has power \(\aleph_1\), each of its finite subsets has a model, and it has no
model \cite[Ch.~II, \S2.2, Example~1, p.~34]{Ebbinghaus1985ExtendedLogics}.
\end{hbtheorem}

With the countable compactness of \(\mathscr L(Q_1)\)
(\cref{thm:p03-keisler-q-one}) this separates countable compactness from
\(\aleph_1\)-compactness, the contrast Barwise draws in
\cite[Ch.~I, \S2.3, p.~16]{Barwise1985BackgroundAims}.

\begin{hbtheorem}{thm:p03-w03-p03-047-sol-interpolation}{proved=annex:p03-w03-p03-047-sol-interpolation}{One-sorted second-order logic has interpolation}
Let \(\mathscr L^2\) be one-sorted second-order logic, with quantifiers over
relations and functions of every finite arity, so that every sentence
contains only finitely many symbols.  If \(\varphi_0\in\mathscr L^2[\tau_0]\),
\(\varphi_1\in\mathscr L^2[\tau_1]\) and \(\varphi_0\models\varphi_1\), then
some \(\psi\in\mathscr L^2[\tau_0\cap\tau_1]\) satisfies
\(\varphi_0\models\psi\) and \(\psi\models\varphi_1\)
\cite[Ch.~II, \S2.1, Warning, p.~33]{Ebbinghaus1985ExtendedLogics}.
\end{hbtheorem}

The proof quantifies away the symbols of \(\varphi_0\) that are not in
\(\tau_1\).  In two-sorted \(\mathscr L^2\) this fails because a sort cannot
be quantified away; Ebbinghaus proves the failure by Craig's method in a
later section and in a later chapter
\cite[Ch.~II, \S2.1, Warning, p.~33]{Ebbinghaus1985ExtendedLogics}.  The
Henkin quantifier of Ebbinghaus's \S2.5 is treated in
\cref{sec:p03-w03-p03-044-po-prefixes}.

\subsection{Countable compactness for regular logics}
\label{sec:p03-general-pinning}

The connection between compactness and well-orderings does not depend on
an embedding into an infinitary language. Relativization and simultaneous
substitution suffice. The latter operation also separates the
interpretations of constants, functions and nullary relations in the
models used in the construction.

\begin{hbconvention}{def:p03-general-pinning-setting}{Scope and cardinality}
In this subsection, \(\mathscr L\) is a one-sorted regular logic in the
sense of \cref{def:p03-w03-p03-047-closure,def:p03-w03-p03-047-relativization}.
Vocabularies may have any set cardinality and may contain constants,
finitary functions and relations, including nullary relations.
Domains are nonempty. Work in classical set theory with Choice.
Use countable compactness and pinning as in
\cref{def:p03-w03-p03-041-compactness,def:p03-w03-p03-041-pinning},
now for these vocabularies and this logic. In particular,
\(w_{\aleph_0}(\mathscr L)\) is the supremum of ordinals pinned by
\emph{at most countably many sentences}; there is no bound on the
size of the vocabulary of those sentences. Pinning requires all models
to give well-orders of the designated relation's field and one model
to give the specified order type, not that every model have that type.
\end{hbconvention}

\begin{hbtheorem}{thm:p03-general-fibres}{proved=annex:p03.general-fibres}{Uniform model fibres}
For any vocabulary \(\tau\), there is an auxiliary vocabulary with a
binary function \(p\) such that, for every \(\varphi\in\mathscr L[\tau]\),
there is a sentence \(\widehat\varphi(d)\) with a distinguished fresh
constant \(d\). Each auxiliary structure \(\mathfrak B\) and element
\(b\in B\) determine a nonempty \(\tau\)-structure \(\mathfrak B_b\)
on
\[
 D_b=\{p^{\mathfrak B}(b,z):z\in B\}
\]
for which
\[
 (\mathfrak B,b)\models_{\mathscr L}\widehat\varphi(d)
 \quad\Longleftrightarrow\quad
 \mathfrak B_b\models_{\mathscr L}\varphi.
\]
The auxiliary vocabulary and the induced structures are the same for
all \(\varphi\). No auxiliary axioms, one for each symbol of \(\tau\),
are required. Any set-indexed family of nonempty \(\tau\)-structures can
be realized, up to isomorphism, as fibres over distinct index elements
in a single auxiliary structure. Its ambient domain may also contain
any prescribed additional set of points, without changing these fibres.
\end{hbtheorem}

\begin{hbtheorem}{thm:p03-general-pinning}{proved=annex:p03.general-pinning}{Countable compactness and well-ordering numbers}
For a one-sorted regular logic \(\mathscr L\), the following are equivalent:
\begin{enumerate}
 \item \(\mathscr L\) is countably compact.
 \item No countable theory of \(\mathscr L\) pins down an infinite ordinal.
 \item \(w_{\aleph_0}(\mathscr L)=\omega\).
\end{enumerate}
If countable compactness fails, countable theories pin down both
\(\omega\) and \(\omega+1\).
This is the regular-logic form of
\cite[Ch.~II, Proposition~5.2.4, p.~62]{Ebbinghaus1985ExtendedLogics}.
\end{hbtheorem}

The general theorem requires no finite-occurrence property, no countable
vocabulary and no prescribed formula grammar. Its proof uses the semantic
closure operations of regularity. The separate
\cref{thm:p03-w03-p03-041-pinning-compactness} remains useful for the
weaker, explicitly syntactic package (F1)--(F4): a sublogic satisfying
that package need not have the full simultaneous-substitution property.


\section{Principal logics: higher order and cardinality quantifiers}
\label{sec:p03-w03-p03-047-principal-logics}

Ebbinghaus's survey of concrete logics serves both directions of the
subject: the general theory is tested on the examples, and the examples
suggest which general properties to study
\cite[Ch.~II, \S2, p.~33]{Ebbinghaus1985ExtendedLogics}.  The examples are
one-sorted except for the higher-order case.

\subsection{Second-order and weak second-order logic}

Second-order logic \(\mathscr L^2\) allows, for each sort, quantification
over \(n\)-ary relations on the domain of that sort
\cite[Ch.~II, Definition~2.1.1, p.~33]{Ebbinghaus1985ExtendedLogics}.  Its
expressive power is bought with the loss of almost every useful
model-theoretic property of first-order logic
(\cref{thm:p03-w03-p03-047-sol-failures}) and with a dependence on set
theory, of which the next theorem is Ebbinghaus's example
\cite[Ch.~II, \S2.1, p.~33]{Ebbinghaus1985ExtendedLogics}.

\begin{hbtheorem}{thm:p03-w03-p03-047-sol-ch}{proved=annex:p03-w03-p03-047-sol-ch}{A second-order sentence valid exactly under the continuum hypothesis}
There is a sentence \(\chi_{\mathrm{CH}}\) of \(\mathscr L^2\) in the empty
vocabulary such that, in \(\mathsf{ZFC}\), \(\chi_{\mathrm{CH}}\) holds in every
structure iff \(2^{\aleph_0}=\aleph_1\).
\end{hbtheorem}

The sentence is built from second-order definitions of ``\((N,<)\) is a copy
of \((\omega,<)\)'', ``\(Y\) is countable'', ``\(Y\) and the domain are
equinumerous'' and ``the domain has power \(2^{\aleph_0}\)''; each is given
and checked in the annex.  Since \(\mathsf{CH}\) is independent of
\(\mathsf{ZFC}\) if \(\mathsf{ZFC}\) is consistent, the validity of
\(\chi_{\mathrm{CH}}\) is not settled by \(\mathsf{ZFC}\); this is the
set-theoretic dependence Ebbinghaus has in mind.

Weak second-order logic \(\mathscr L^{w2}\) has the same syntax as
\(\mathscr L^2\), but its relation variables range only over finite
relations \cite[Ch.~II, Definition~2.1.2, p.~33]{Ebbinghaus1985ExtendedLogics}.
In the one-sorted case treated here, \(\exists X\) with \(X\) an \(n\)-ary
relation variable means ``for some finite \(n\)-ary relation on the
domain''.

\begin{hbtheorem}{thm:p03-w03-p03-047-weak-sol}{proved=annex:p03-w03-p03-047-weak-sol}{Properties of weak second-order logic}
Let \(\mathscr L^{w2}\) be one-sorted weak second-order logic.
\begin{enumerate}
 \item The sentence \(\exists X\,\forall x\,Xx\) holds in a structure iff
       its domain is finite.  More generally, for every formula
       \(\varphi(x)\) of \(\mathscr L^{w2}\), with \(X\) a unary variable not
       occurring in \(\varphi\), the formula
       \(\neg\exists X\,\forall x\,(\varphi(x)\to Xx)\) holds iff infinitely
       many elements satisfy \(\varphi\).  Hence every
       \(\mathscr L(Q_0)\)-formula has an \(\mathscr L^{w2}\)-equivalent,
       obtained by a computable translation.
 \item In the vocabulary \(\{\cdot,{}^{-1},e\}\) a group satisfies
       \(\forall x\,\exists X\,(Xe\land\forall y\,(Xy\to X(y\cdot x)))\) iff
       it is a torsion group.
 \item Some \(\mathscr L^{w2}\)-sentence of vocabulary \(\Lang_A\cup\{<\}\)
       has exactly the structures isomorphic to \((\mathcal N,<)\) as models.
 \item \(\mathscr L^{w2}\) is not \(\aleph_0\)-compact, and, under a coding
       in which the translation of part~1 is computable, the valid
       \(\mathscr L^{w2}\)-sentences of \(\Lang_A\cup\{<\}\) do not form a
       recursively enumerable set.
 \item For every \(\mathscr L^{w2}\)-formula \(\varphi\) without free
       relation variables there is an \(\mathscr L_{\omega_1\omega}\)-formula
       \(\varphi^\ast\) with the same free individual variables and the same
       satisfying assignments in every structure.  Consequently
       \(\mathscr L^{w2}\) has the L\"owenheim--Skolem property of
       \cref{def:p03-w03-p03-047-ebbinghaus-properties}.
\end{enumerate}
\end{hbtheorem}

Ebbinghaus adds that, since arithmetical \(\mathscr L^{w2}\)-truth is
implicitly definable, the method of his \S7.3 shows that the Beth property,
and with it the interpolation property, fail for \(\mathscr L^{w2}\)
\cite[Ch.~II, \S2.1, p.~34]{Ebbinghaus1985ExtendedLogics}.  The failure of
Beth's theorem is proved, with true arithmetic as the implicitly but not
explicitly definable relation, for Barwise's formulation of weak
second-order logic over hereditarily finite sets in
\cref{thm:p02-w02c-p02-003-wso-beth-failure}.  For the finite-relation
formulation used here the failure is part~2 of
\cref{thm:p03-w03-p03-047-mostowski-general}: \(\mathscr L^{w2}\) defines its
own satisfaction relation over \((\mathcal N,<)\), which is then implicitly
but not explicitly definable.  The interpolation property fails too, by
\cref{thm:p03-w03-p03-047-craig-beth}, which needs only the atom, negation
and conjunction properties.

\subsection{Cardinality quantifiers}

If \(\mathscr L_{\omega\omega}\) is extended by one unary quantifier that is
monotone in the sense of Ebbinghaus's Theorem~4.2.3, the extension is regular
exactly when the quantifier is some \(Q_\alpha\); Ebbinghaus cites this from
a later chapter of the volume \cite[Ch.~II, \S2.2, p.~34]{Ebbinghaus1985ExtendedLogics},
and it is not proved here.  The logic in focus is
\(\mathscr L_{\omega\omega}(Q_1)\) of \cref{def:p03-w03-p03-047-mostowski-q}.
It is \(\aleph_0\)-compact, by Fuhrken's theorem, and recursively enumerable
for consequence, by Vaught's; Keisler gave a complete calculus by adding four
axiom schemata to a first-order calculus
\cite[Ch.~II, \S2.2, Example~1, p.~34]{Ebbinghaus1985ExtendedLogics}.  These
schemata are those of \cref{def:p02-w02c-p02-003-kq-calculus}, where the
completeness and compactness theorems are proved
(\cref{thm:p02-w02c-p02-003-keisler-completeness,thm:p02-w02c-p02-003-keisler-compactness}).
Ebbinghaus locates the expressive power of
\(\mathscr L_{\omega\omega}(Q_1)\) beyond first-order logic in the
characterization of \emph{\(\aleph_1\)-like orderings}, the linear orders
\(\mathfrak A=(A,<^{\mathfrak A})\) satisfying
\(Q_1x\,x=x\land\forall y\,\neg Q_1x\,x<y\)
\cite[Ch.~II, \S2.2, Example~1, p.~34]{Ebbinghaus1985ExtendedLogics}.

\begin{hbtheorem}{thm:p03-w03-p03-047-q1-ls-consequence}{proved=annex:p03-w03-p03-047-q1-ls-consequence}{L\"owenheim--Skolem down to \(\aleph_1\) and enumerability of consequence for \(\mathscr L_{\omega\omega}(Q_1)\)}
\begin{enumerate}
 \item Let \(\Phi\) be a set of at most \(\aleph_1\) sentences of
       \(\mathscr L_{\omega\omega}(Q_1)\).  If \(\Phi\) has a model, it has a
       model of power at most \(\aleph_1\).
 \item For a countable vocabulary as in part~1 of
       \cref{thm:p03-w03-p03-047-q1-mostowski-facts} and a recursively
       enumerable set \(\Sigma\) of sentences, the set of sentences
       \(\varphi\) with \(\Sigma\models\varphi\) is recursively enumerable.
\end{enumerate}
\end{hbtheorem}

Part~1 cannot be lowered to \(\aleph_0\): \(Q_1x\,(x=x)\) has models and none
of power \(\aleph_0\) (part~4 of \cref{thm:p02-w02-p02-005-bw-extended-ls}).
Compactness fails at \(\aleph_1\) (\cref{thm:p03-w03-p03-047-q1-not-aleph1}).
Ebbinghaus also reports an omitting types theorem for
\(\mathscr L_{\omega\omega}(Q_1)\), the lack of satisfactory preservation
theorems, failure of interpolation (Keisler), and failure of the Beth
property (H.~Friedman)
\cite[Ch.~II, \S2.2, Example~1, p.~34]{Ebbinghaus1985ExtendedLogics}.
The interpolation failure is proved in
\cref{thm:p03-lq-qone-interpolation}. The other assertions in this
list are reported here without proofs.

\begin{hbexample}{def:p03-w03-p03-047-keisler-entailment}{Keisler's entailment}
Let \(E\) be binary and \(R,S\) unary.  Let \(\varphi(E)\) say that \(E\) is
an equivalence relation all of whose classes are uncountable:
\(E\) is reflexive, symmetric and transitive, and \(\forall x\,Q_1y\,Exy\).
Let
\[
 \varphi_0(E,R)=\varphi(E)\land\neg Q_1x\,Rx\land\mathrm{Rep}(E,R),\qquad
 \varphi_1(E,S)=\varphi(E)\land Q_1x\,Sx\land\mathrm{Rep}(E,S),
\]
where \(\mathrm{Rep}(E,P)\) is
\(\forall x\,\exists y\,(Py\land Exy)\land\forall x\,\forall y\,(Px\land Py\land Exy\to x=y)\),
saying that \(P\) is a set of representatives.  So \(\varphi_0\) says that
\(R\) is a countable set of representatives and \(\varphi_1\) that \(S\) is
an uncountable one, and
\[
 (\ast)\qquad \varphi_0(E,R)\models\neg\varphi_1(E,S)
\]
\cite[Ch.~II, \S2.2, p.~35]{Ebbinghaus1985ExtendedLogics}.  The entailment
holds because two sets of representatives of one equivalence relation are in
bijection (send each element of \(R\) to the element of \(S\) in its class),
so they cannot be countable and uncountable at once.  The common vocabulary
is \(\{E\}\).  Ebbinghaus states that \((\ast)\) has no
\(\mathscr L_{\omega\omega}(Q_1)\)-interpolant, referring to his \S4.2.
\Cref{thm:p03-lq-qone-interpolation} proves this by an explicit
unbounded comparison system between structures with countably and
uncountably many uncountable classes.
\end{hbexample}

The next two logics remove this counterexample.  Magidor and Malitz's
quantifiers \(Q_1^n\), \(n\ge1\), are \(n\)-ary:
\(\mathfrak A\models Q_1^n\mathbf x\,\varphi(\mathbf x)\) iff some uncountable
\(M\subseteq A\) has \(\mathfrak A\models\varphi[\mathbf b]\) for all
\(\mathbf b\in M^n\) \cite[Ch.~II, \S2.2, Example~2, p.~35]{Ebbinghaus1985ExtendedLogics}.
Ebbinghaus observes that \(Q_1x\,\varphi(x)\) has the same meaning as
\(\exists\) uncountable \(X\,\forall x\,(\neg Xx\lor\varphi(x))\) and as
\(\neg\exists\) countable \(X\,\forall x\,(Xx\lor\neg\varphi(x))\), and
considers logics that quantify over uncountable, or over countable, subsets.
A set variable \(X\) occurs \emph{negatively} (\emph{positively}) in \(\varphi\),
with \(\neg,\land,\lor\) the only connectives, if some occurrence of \(X\)
lies in the scope of an odd (even) number of negations.  The logic
\(\mathscr L_{\omega\omega}(\mathrm{pos})\) (\(\mathscr L_{\omega\omega}(\mathrm{neg})\))
allows \(\exists X\varphi\), with \(X\) ranging over countable (uncountable)
subsets, only when \(X\) occurs at most positively (negatively) in \(\varphi\)
\cite[Ch.~II, \S2.2, Example~3, p.~36]{Ebbinghaus1985ExtendedLogics}.

\begin{hbtheorem}{thm:p03-w03-p03-047-q1-interpolants}{proved=annex:p03-w03-p03-047-q1-interpolants}{Interpolants for Keisler's entailment in stronger logics}
In the notation of \cref{def:p03-w03-p03-047-keisler-entailment}:
\begin{enumerate}
 \item \(Q_1x\,\psi(x)\) and \(Q_1^2xy\,(\psi(x)\land\psi(y))\) are
       equivalent for every formula \(\psi(x)\), so \(\varphi(E)\) has an
       equivalent in \(\mathscr L_{\omega\omega}(Q_1^2)\); and
       \(\varphi(E)\land\neg Q_1^2xy\,(x=y\lor\neg Exy)\) is an interpolant
       for \((\ast)\) in \(\mathscr L_{\omega\omega}(Q_1^2)\).
 \item \(\varphi(E)\land\exists X\,\forall y\,\exists x\,(Xx\land Exy)\), with
       \(X\) ranging over countable subsets and \(Q_1\) read by the countable
       form above, is a sentence of \(\mathscr L_{\omega\omega}(\mathrm{pos})\)
       and an interpolant for \((\ast)\).
\end{enumerate}
Given \(\varphi(E)\), both interpolants say that \(E\) has only countably many
classes.
\end{hbtheorem}

By part~1, \(\mathscr L_{\omega\omega}(Q_1^2)\) overcomes Keisler's
counterexample, so, given the non-existence claim of
\cref{def:p03-w03-p03-047-keisler-entailment}, it is strictly stronger than
\(\mathscr L_{\omega\omega}(Q_1)\)
\cite[Ch.~II, \S2.2, Example~2, p.~35]{Ebbinghaus1985ExtendedLogics}.  The
remaining claims of Example~2 are reported and not proved here: under
\(V=L\), or under \(\diamondsuit_{\aleph_1}\), Magidor and Malitz showed
\(\mathscr L_{\omega\omega}(Q_1^n\mid n\ge1)\) to be \(\aleph_0\)-compact; by
a result of Shelah it is consistent that \(\mathscr L_{\omega\omega}(Q_1^2)\)
is not \(\aleph_0\)-compact; Suslin trees are characterizable in
\(\mathscr L_{\omega\omega}(Q_1^2)\); no \(\mathscr L_{\omega\omega}(Q_1^n)\)
has the Beth property (Badger); and \(\mathscr L_{\omega\omega}(Q_1^{n+1})\)
is more expressive than \(\mathscr L_{\omega\omega}(Q_1^n)\) (Garavaglia and
Shelah) \cite[Ch.~II, \S2.2, Example~2, p.~35]{Ebbinghaus1985ExtendedLogics}.

\begin{hbtheorem}{thm:p03-w03-p03-047-set-quantifier-failures}{proved=annex:p03-w03-p03-047-set-quantifier-failures}{Quantifying over all uncountable or all countable subsets}
\begin{enumerate}
 \item A linear ordering is isomorphic to \((\omega_1,<)\) iff it is
       \(\aleph_1\)-like and each of its uncountable subsets has a least
       element.  Hence, with unrestricted quantification over uncountable
       subsets, some sentence has exactly the copies of \((\omega_1,<)\) as
       models, and the logic is not \(\aleph_0\)-compact.
 \item With unrestricted quantification over countable subsets, some
       sentence has exactly the copies of \((\omega,<)\) as models, and the
       logic is not \(\aleph_0\)-compact.
\end{enumerate}
\end{hbtheorem}

This is why Ebbinghaus restricts the polarity of the quantified variable
\cite[Ch.~II, \S2.2, Example~3, p.~36]{Ebbinghaus1985ExtendedLogics}.
He reports that \(\mathscr L_{\omega\omega}(\mathrm{neg})\) extends
\(\mathscr L_{\omega\omega}(Q_1^n\mid n\ge1)\) but still characterizes
\((\omega_1,<)\) (Stavi), and that \(\mathscr L_{\omega\omega}(\mathrm{pos})\)
is \(\aleph_0\)-compact and recursively enumerable for consequence
\cite[Ch.~II, \S2.2, Example~3, p.~36]{Ebbinghaus1985ExtendedLogics}; these
are not proved here.  The property of \(\mathscr L_{\omega\omega}(\mathrm{pos})\)
that Ebbinghaus obtains by an induction on formulas is the following; the
polarity hypothesis on the free variable is stated explicitly.

\begin{hbtheorem}{thm:p03-w03-p03-047-pos-monotone}{proved=annex:p03-w03-p03-047-pos-monotone}{Positive countable-set quantification is upward monotone}
Let \(\varphi\) be an \(\mathscr L_{\omega\omega}(\mathrm{pos})\)-formula
built with \(\neg,\land,\lor\), the quantifiers \(\exists,\forall\) over
elements and \(\exists\) over countable sets, and let \(X\) be a set variable
free in \(\varphi\).  If \(X\) occurs only positively in \(\varphi\), then for
all countable \(M\subseteq M'\subseteq A\) and all values of the other free
variables, \(\mathfrak A\models\varphi[M,\ldots]\) implies
\(\mathfrak A\models\varphi[M',\ldots]\).  If \(X\) occurs only negatively,
\(\mathfrak A\models\varphi[M',\ldots]\) implies
\(\mathfrak A\models\varphi[M,\ldots]\)
\cite[Ch.~II, \S2.2, Example~3, p.~36]{Ebbinghaus1985ExtendedLogics}.
\end{hbtheorem}

\begin{hbwarning}{def:p03-w03-p03-047-pos-polarity}{The polarity hypothesis cannot be dropped}
For the formula \(\neg Xx\), in which \(X\) occurs negatively, take
\(A=\omega\), \(M=\{0\}\), \(M'=\{0,1\}\) and \(x\) assigned \(1\): the
formula holds at \(M\) and fails at \(M'\).  The monotonicity of
\cref{thm:p03-w03-p03-047-pos-monotone} is a property of positive
occurrences, which the syntax of \(\mathscr L_{\omega\omega}(\mathrm{pos})\)
guarantees for every bound set variable at its quantifier.  The logic can
assert that some countable set is large enough for a purpose, but it cannot
say that a countable set is small.
\end{hbwarning}

\section{Comparing logics}
\label{sec:p03-w03-p03-047-comparing}

Ebbinghaus turns the informal comparisons of expressive power into
definitions by means of elementary and projective classes
\cite[Ch.~II, \S3, p.~41]{Ebbinghaus1985ExtendedLogics}.  Throughout,
\(\operatorname{Mod}_{\mathscr L}(\varphi)\) is the class of models of an
\(\mathscr L\)-sentence \(\varphi\) and \(\overline K=\operatorname{Str}[\tau]\setminus K\)
for a class \(K\) of \(\tau\)-structures.

\begin{hbdefinition}{def:p03-w03-p03-047-ec-pc-rpc}{Elementary, projective and relativized projective classes (Ebbinghaus, Definition~3.1.1)}
Let \(\mathscr L\) be a logic and \(K\) a class of \(\tau\)-structures.
\begin{enumerate}
 \item \(K\) is \emph{EC in \(\mathscr L\)}, \(K\in\mathrm{EC}_{\mathscr L}\),
       if \(K=\operatorname{Mod}_{\mathscr L}(\varphi)\) for some
       \(\varphi\in\mathscr L[\tau]\).
 \item \(K\) is \emph{PC in \(\mathscr L\)}, \(K\in\mathrm{PC}_{\mathscr L}\),
       if there are \(\tau'\supseteq\tau\), with the same sort symbols as
       \(\tau\) in the many-sorted case, and \(K'\in\mathrm{EC}_{\mathscr L}\)
       of \(\tau'\)-structures with
       \(K=\{\mathfrak A\restriction\tau:\mathfrak A\in K'\}\).
 \item \(K\) is \emph{RPC in \(\mathscr L\)}, \(K\in\mathrm{RPC}_{\mathscr L}\),
       if in the one-sorted case there are \(\tau'\supseteq\tau\), a unary
       \(U\in\tau'\setminus\tau\) and \(K'\in\mathrm{EC}_{\mathscr L}\) of
       \(\tau'\)-structures with
       \[
        K=\{(\mathfrak A\restriction\tau)|U^{\mathfrak A}:
          \mathfrak A\in K'\text{ and }U^{\mathfrak A}\text{ is }\tau\text{-closed in }\mathfrak A\},
       \]
       and in the many-sorted case there are \(\tau'\supseteq\tau\), now
       possibly with new sort symbols, and \(K'\in\mathrm{EC}_{\mathscr L}\)
       with \(K=\{\mathfrak A\restriction\tau:\mathfrak A\in K'\}\).
\end{enumerate}
\cite[Ch.~II, Definition~3.1.1, pp.~41--42]{Ebbinghaus1985ExtendedLogics}.
Here \((\mathfrak A\restriction\tau)|U^{\mathfrak A}\) is the substructure
of \(\mathfrak A\restriction\tau\) with domain \(U^{\mathfrak A}\), which
exists because \(U^{\mathfrak A}\) is closed under the functions and
contains the constants of \(\tau\).  The PC classes of
\cref{def:p03-w03-p03-047-pc-delta} are the one-sorted PC classes of this
definition.
\end{hbdefinition}

Ebbinghaus remarks that for the usual logics, and for every regular logic
over finite vocabularies, a class of one-sorted structures is RPC in the
one-sorted sense iff it is RPC in the many-sorted sense, by unifying
domains in one direction and splitting them in the other; that PC is
included in RPC for every logic containing sentences such as
\(\forall x\,Ux\), strictly for \(\mathscr L_{\omega\omega}\) and not
strictly for \(\mathscr L_{\omega_1\omega}\); and that a PC class need not be
EC, even in first-order logic \cite[Ch.~II, Remarks~3.1.2, p.~42]{Ebbinghaus1985ExtendedLogics}.
The inclusion and the first-order counterexample are proved next; the
equivalence of the two RPC notions and the two strictness claims are
reported, not proved, here.

\begin{hbtheorem}{thm:p03-w03-p03-047-pc-basic}{proved=annex:p03-w03-p03-047-pc-basic}{PC inside RPC; infinity is PC but not EC}
\begin{enumerate}
 \item Let \(\mathscr L\) be one-sorted with the conjunction property, and
       suppose that for every \(\tau'\) and unary \(U\in\tau'\) some
       \(\mathscr L[\tau']\)-sentence has as models exactly the
       \(\tau'\)-structures with \(U^{\mathfrak A}=A\).  Then
       \(\mathrm{PC}_{\mathscr L}\subseteq\mathrm{RPC}_{\mathscr L}\).
 \item The class of infinite sets (structures of the empty vocabulary) is PC
       in \(\mathscr L_{\omega\omega}\), witnessed by the sentence of
       vocabulary \(\{f\}\), \(f\) unary,
       \(\forall x\,\forall y\,(fx=fy\to x=y)\land\exists y\,\forall x\,\neg fx=y\);
       it is not EC in \(\mathscr L_{\omega\omega}\).
\end{enumerate}
\end{hbtheorem}

\cref{thm:p03-w03-p03-047-craig-separation} part~1 is the one-sorted case
of Ebbinghaus's Proposition~3.1.3, that for a logic with the negation
property interpolation is equivalent to the separation of disjoint PC classes
(RPC classes in the many-sorted case) by an EC class
\cite[Ch.~II, Proposition~3.1.3, p.~42]{Ebbinghaus1985ExtendedLogics}.  The
question whether a class \(K\) with \(K\) and \(\overline K\) both (R)PC is
EC leads to \(\Delta\)-interpolation, of which interpolation is in this sense
a generalization.

\begin{hbdefinition}{def:p03-w03-p03-047-strength}{Strength of logics (Ebbinghaus, Definitions~3.1.4--3.1.6)}
Let \(\mathscr L,\mathscr L^*\) be logics.
\begin{enumerate}
 \item \(\mathscr L\le\mathscr L^*\) (\(\mathscr L^*\) is \emph{as strong as}
       \(\mathscr L\)) if every class EC in \(\mathscr L\) is EC in
       \(\mathscr L^*\); \(\mathscr L\equiv\mathscr L^*\) (\emph{equally
       strong}) if \(\mathscr L\le\mathscr L^*\) and \(\mathscr L^*\le\mathscr L\);
       \(\mathscr L<\mathscr L^*\) (\emph{stronger}) if \(\mathscr L\le\mathscr L^*\)
       and not \(\mathscr L^*\le\mathscr L\).
 \item \(\mathscr L\le_{\equiv}\mathscr L^*\) if for all \(\tau\) and all
       \(\mathfrak A,\mathfrak B\in\operatorname{Str}[\tau]\),
       \(\mathfrak A\equiv_{\mathscr L^*}\mathfrak B\) implies
       \(\mathfrak A\equiv_{\mathscr L}\mathfrak B\), where
       \(\equiv_{\mathscr L}\) means satisfying the same sentences of
       \(\mathscr L[\tau]\).
 \item \(\mathscr L\le_{\mathrm{PC}}\mathscr L^*\)
       (\(\mathscr L\le_{\mathrm{RPC}}\mathscr L^*\)) if every class PC
       (RPC) in \(\mathscr L\) is PC (RPC) in \(\mathscr L^*\); the relations
       \(<_{(\mathrm R)\mathrm{PC}}\) and \(\equiv_{(\mathrm R)\mathrm{PC}}\)
       are defined from these as in~1.
\end{enumerate}
\cite[Ch.~II, Definitions~3.1.4--3.1.6, p.~43]{Ebbinghaus1985ExtendedLogics}.
\end{hbdefinition}

Definition~1 identifies the meaning of a sentence with its model class and
asks that every \(\mathscr L\)-sentence have an \(\mathscr L^*\)-sentence
with the same meaning; definition~2 asks instead that structures told apart
by \(\mathscr L\) be told apart by \(\mathscr L^*\)
\cite[Ch.~II, \S3.1, pp.~42--43]{Ebbinghaus1985ExtendedLogics}.

\begin{hbtheorem}{thm:p03-w03-p03-047-strength-basic}{proved=annex:p03-w03-p03-047-strength-basic}{Basic facts on comparison}
\begin{enumerate}
 \item \(\le\) is reflexive and transitive, so \(\equiv\) is an equivalence
       relation and \(\le\) is a partial ordering of the
       \(\equiv\)-classes; the same holds for \(\le_{\mathrm{PC}}\) and
       \(\le_{\mathrm{RPC}}\).
 \item If \(\mathscr L\le\mathscr L^*\) then
       \(\mathscr L\le_{\equiv}\mathscr L^*\).
\end{enumerate}
\end{hbtheorem}

The converse of part~2 can fail: Ebbinghaus notes that
\(\mathscr L_{\infty G}\le_{\equiv}\mathscr L_{\infty\omega}\), because
\(\mathscr L_{\infty G}\) has the Karp property, while
\(\mathscr L_{\infty\omega}<\mathscr L_{\infty G}\)
\cite[Ch.~II, \S3.1, p.~43]{Ebbinghaus1985ExtendedLogics}; the example rests
on the game logic of a later chapter and is reported, not proved, here.
From the examples of the preceding section he reads off the strict chains
\[
 \mathscr L_{\omega\omega}<\mathscr L^{w2}<\mathscr L^2,\qquad
 \mathscr L_{\omega\omega}<\mathscr L_{\omega\omega}(Q_1)<\mathscr L_{\omega\omega}(Q_1^2)<\cdots,
\]
\[
 \mathscr L_{\omega\omega}(Q_1)<\mathscr L_{\omega\omega}(\mathrm{pos})<\mathscr L_{\omega\omega}(\mathrm{aa}),\qquad
 \mathscr L_{\omega\omega}<\mathscr L_{\omega\omega}(I)<\mathscr L_{\omega\omega}(Q^{\mathrm R}),
\]
together with \(\mathscr L^{w2}<\mathscr L_{\omega_1\omega}\),
\(\mathscr L^2\not\le\mathscr L_{\omega_1\omega}\) and
\(\mathscr L_{\omega_1\omega}\not\le\mathscr L^2\); and for the class
\(\mathfrak K\) of \(\aleph_1\)-like orderings,
\(\mathscr L_{\omega\omega}(\mathfrak K)\le\mathscr L_{\omega\omega}(Q_1)\)
while the sentence \(Q_1x\,x=x\) shows that the converse fails.  That
failure is why the projective comparisons are introduced
\cite[Ch.~II, \S3.1, p.~43]{Ebbinghaus1985ExtendedLogics}.  These
comparisons are reported, not proved, here.

\begin{hbtheorem}{thm:p03-w03-p03-047-rpc-descent}{proved=annex:p03-w03-p03-047-rpc-descent}{Compactness and L\"owenheim--Skolem descend along \(\le_{\mathrm{RPC}}\) (Ebbinghaus, Proposition~3.1.9)}
Let \(\mathscr L\le_{\mathrm{RPC}}\mathscr L^*\) in the many-sorted reading
of \cref{def:p03-w03-p03-047-ec-pc-rpc}, and let \(\kappa\) be infinite.
\begin{enumerate}
 \item If \(\mathscr L^*\) is \(\kappa\)-compact, so is \(\mathscr L\);
       hence if \(\mathscr L^*\) is compact, so is \(\mathscr L\).
 \item If \(\mathscr L^*\) has the L\"owenheim--Skolem property down to
       \(\kappa\), so does \(\mathscr L\).
\end{enumerate}
\cite[Ch.~II, Proposition~3.1.9, p.~44]{Ebbinghaus1985ExtendedLogics}.
\end{hbtheorem}

\begin{hbmethod}{def:p03-w03-p03-047-reduction-scheme}{The reduction method}
To show that a logic \(\mathscr L\) has a property \(P\): (a) find a logic
\(\mathscr L^*\) with \(\mathscr L\le_{\mathrm{RPC}}\mathscr L^*\); (b) prove
\(P\) for \(\mathscr L^*\); (c) verify that \(P\) passes down from
\(\mathscr L^*\) to \(\mathscr L\).  For \(\kappa\)-compactness and the
L\"owenheim--Skolem property down to \(\kappa\), step~(c) is supplied once
and for all by \cref{thm:p03-w03-p03-047-rpc-descent}
\cite[Ch.~II, \S3.1, pp.~44--45]{Ebbinghaus1985ExtendedLogics}.  Often
\(\mathscr L^*\) is first-order logic restricted to a class of structures,
such as \(\aleph_1\)-like orderings.
\end{hbmethod}

The source's own instance, Proposition~3.1.7
(\(\mathscr L_{\omega\omega}(Q_1)\le_{\mathrm{RPC}}\mathscr L_{\omega\omega}(\mathfrak K)\)
for \(\aleph_1\)-like orderings) and Corollary~3.1.8 (countable compactness
of \(\mathscr L_{\omega\omega}(Q_1)\)), rests on Theorem~2.6.1 of the
preceding section; countable compactness of \(\mathscr L(Q_1)\) by another
route is \cref{thm:p03-keisler-q-one}.

\section{Lindstr\"om quantifiers and back-and-forth}
\label{sec:p03-lindstrom-quantifiers}

A generalized quantifier tests a structure defined by its argument
formulas. The defining class is part of the quantifier's semantics;
it need not be first-order definable. This description separates three
questions: what class the quantifier recognizes, which closure operations
its logic retains, and which partial isomorphisms preserve its formulas.

\begin{hbconvention}{def:p03-lq-setting}{Structures and formula parameters}
Work with one-sorted structures with nonempty domains, in classical set
theory with Choice. All function and relation arities are finite.
The first-order and \(L_{\infty\omega}\) languages used below have
finitely many free variables in each formula. For abstract logics,
expressive equivalence concerns sentences over the same vocabulary.
The translations also preserve formulas with finitely many free
parameters, represented by fresh constants as in
\cref{def:p03-w03-p03-047-closure}. No claim is made that a formula
with infinitely many free parameters has a finite-parameter translation.
\end{hbconvention}

\subsection{Quantifiers from classes of structures}

\begin{hbdefinition}{def:p03-lq-full}{A full-domain class quantifier}
Let \(\sigma\) be a finite vocabulary, called the \emph{pattern}, and
let \(K\) be an isomorphism-closed class of \(\sigma\)-structures.
For each relation \(R\) of arity \(r\), an application of \(Q_K\)
has an operand \(\alpha_R(\bar x_R,\bar u)\) with \(r\) bound
variables. For each function \(f\) of arity \(r\), it has a graph
operand \(\beta_f(\bar x_f,y_f,\bar u)\), and for each constant
\(c\), a singleton operand \(\gamma_c(z_c,\bar u)\).
Each bound list consists of distinct variables; these lists are
chosen disjoint from the free parameter list \(\bar u\).
Write the application as
\[
 Q_K\bigl((\bar x_R\,\alpha_R)_R,
          (\bar x_f y_f\,\beta_f)_f,(z_c\,\gamma_c)_c\bigr).
\]
Each list binds variables only in its own operand.
In a structure \(\mathfrak A\), with parameters \(\bar a\), it is
true precisely when the operands define, on the whole domain \(A\),
the relations, total function graphs and unique constants of a
\(\sigma\)-structure belonging to \(K\).
If a function operand is not a total single-valued graph, or a
constant operand is not a singleton, the application is false.
A nullary relation has an operand with no bound variables.

For a family of such quantifiers, \(\operatorname{FO}(Q_i:i\in J)\)
is the ordinary finite formula language obtained from first-order
atoms by Boolean operations, first-order quantifiers and these
applications. Each formula uses only finitely many members of the
family. With a general base logic, \(\mathscr L(Q_i:i\in J)\)
is the analogous finite closure starting from its formulas.
This is the class-quantifier construction of
\cite[Ch.~II, Definition~4.1.1, pp.~49--50]{Ebbinghaus1985ExtendedLogics}.
\end{hbdefinition}

\begin{hbexample}{def:p03-lq-examples}{Standard quantifiers as class tests}
The following choices of \(K\) recover familiar operations directly
from \cref{def:p03-lq-full}.
\begin{enumerate}
 \item For a unary predicate \(P\), take the structures with
       \(P\ne\varnothing\). Then \(Q_Kx\,\alpha\) is
       \(\exists x\,\alpha\).
 \item For an \(r\)-ary predicate \(R\), with \(r\geq1\), take the structures
       containing \(C\subseteq A\) with \(|C|\geq\aleph_\alpha\)
       and \(C^r\subseteq R\). The result is the homogeneous-set
       quantifier \(Q^r_\alpha\); at \(r=1\) it is the usual
       cardinality quantifier.
 \item For binary \(R\), take the strict linear orders of the whole
       domain with cofinality \(\omega\). This is
       \(Q^{\mathrm{cf}\,\omega}\). Taking all well-orders of the
       whole domain instead gives \(Q^{\mathrm{wo}}\).
 \item For unary \(P\), take \(|P|=|A|\). This is Chang's
       full-domain quantifier \(Q^{\mathrm C}\).
\end{enumerate}
In the order examples, an operand defining a relation that is not a
linear order is rejected. The last example compares with the whole
ambient domain, not a second supplied predicate
\cite[Ch.~II, Example~4.1.2, p.~50]{Ebbinghaus1985ExtendedLogics}.
\end{hbexample}

\begin{hbtheorem}{thm:p03-lq-representation}{proved=annex:p03.lq.1}{Representation by class quantifiers}
Let \(\mathscr L\) be a one-sorted regular logic.
\begin{enumerate}
 \item Every isomorphism-closed class \(K\) of a finite pattern
       vocabulary is defined by one sentence of
       \(\operatorname{FO}(Q_K)\).
 \item If \(K\) is already defined by a sentence of \(\mathscr L\),
       every formula of \(\mathscr L(Q_K)\) has an equivalent
       \(\mathscr L\)-formula.
 \item Suppose \(\mathscr L\) is \emph{finitary}: each formula
       belongs to the logic of a finite sub-vocabulary, with its
       finite list of parameter constants retained. Then
       \(\mathscr L\) is expressively equivalent to first-order
       logic with the quantifiers \(Q_K\) for all its definable
       classes of finite vocabulary.
\end{enumerate}
The family in the last clause need not be finite
\cite[Ch.~II, Theorem~4.1.3, p.~51]{Ebbinghaus1985ExtendedLogics}.
\end{hbtheorem}

\begin{hbtheorem}{thm:p03-lq-second-order}{proved=annex:p03.lq.2}{Second-order logic needs an infinite family}
Full second-order logic is not expressively equivalent to
\(\operatorname{FO}(Q_1,\ldots,Q_s)\) for any finite family of
ordinary Lindstr\"om quantifiers of finite pattern vocabulary.
Consequently the representation of second-order logic in
\cref{thm:p03-lq-representation} cannot use a finite family.
This is the non-finite-generation assertion following
\cite[Ch.~II, Theorem~4.1.3, p.~51]{Ebbinghaus1985ExtendedLogics}.
\end{hbtheorem}

\subsection{Relativized domains}

\begin{hbdefinition}{def:p03-lq-relative}{A domain-relativized class quantifier}
An application of \(Q_K^*\) has, in addition to the operands of
\cref{def:p03-lq-full}, a domain operand \(\delta(v,\bar u)\).
In \(\mathfrak A\) its domain is
\(D=\{b\in A:\mathfrak A\models\delta(b,\bar a)\}\).
The application is true exactly when \(D\ne\varnothing\) and the
other operands, evaluated in \(\mathfrak A\) but restricted to tuples
from \(D\), define a \(\sigma\)-structure in \(K\).
Function operands must be total single-valued graphs on \(D\), and
constant operands must define singletons in \(D\).
No condition is imposed on those operands outside \(D\).
This is the domain variant of
\cite[Ch.~II, Definition~4.1.4, p.~51]{Ebbinghaus1985ExtendedLogics}.
\end{hbdefinition}

\begin{hbtheorem}{thm:p03-lq-domain}{proved=annex:p03.lq.3}{Domain translation and regularity}
Let \(K\) have finite pattern \(\sigma\), and let \(U\) be a fresh
unary relation. Define \(K^U\) to consist of the
\((\sigma\cup\{U\})\)-structures in which \(U\) is nonempty and
\(\sigma\)-closed and the induced \(\sigma\)-structure on \(U\)
belongs to \(K\). Then:
\begin{enumerate}
 \item \(Q_K\) is expressed by \(Q_K^*\) with domain formula \(v=v\).
 \item The logics \(\operatorname{FO}(Q_K^*)\) and
       \(\operatorname{FO}(Q_{K^U})\) are expressively equivalent.
       The translations preserve free parameters and reject malformed
       function graphs and constant descriptions correctly.
 \item For any family of finite patterns,
       \(\operatorname{FO}(Q_i^*:i\in J)\) is regular, with
       substitution interpreted on admissible interpretations as in
       \cref{def:p03-w03-p03-047-relativization}.
\end{enumerate}
The class conversion is
\cite[Ch.~II, Proposition~4.1.5, p.~52]{Ebbinghaus1985ExtendedLogics}.
\end{hbtheorem}

\begin{hbcorollary}{thm:p03-lq-minimal}{proved=annex:p03.lq.4}{The regular extension generated by a class}
Let \(K\) have finite pattern vocabulary. Every regular logic in
which \(K\) is definable contains the expressive power of
\(\operatorname{FO}(Q_K^*)\). Thus this logic is the least regular
extension of first-order logic defining \(K\), up to expressive
equivalence. More generally, if \(\mathscr L(Q_K^*)\) is itself
regular, it is the least regular extension of \(\mathscr L\)
defining \(K\).
\end{hbcorollary}

\begin{hbtheorem}{thm:p03-lq-chang-boundary}{proved=annex:p03.lq.5}{A failure of relativization}
The logic \(\operatorname{FO}(Q^{\mathrm C})\) is not closed under
relativization to a nonempty unary domain. There are structures in the
vocabulary \(\{U,V\}\) satisfying the same sentences of this logic,
but differing on whether \(|U\cap V|=|U|\).
Thus ordinary full-domain class quantifiers do not automatically give
regular extensions, even when their defining classes are monotone.
\end{hbtheorem}

\subsection{A standard part is not a quantifier}

\begin{hbdefinition}{def:p03-lq-standard-part}{Standard-part semantics}
Fix a finite pattern \(\sigma\), a fresh unary \(U\), and an
isomorphism-closed class \(K\) of \(\sigma\)-structures.
The standard-part semantics \(\operatorname{FO}[K]\) uses ordinary
first-order sentences over vocabularies containing \(\sigma\cup\{U\}\),
but permits only structures in which \(U\) is nonempty and
\(\sigma\)-closed and its induced \(\sigma\)-structure belongs to \(K\).
On permitted structures satisfaction is ordinary first-order satisfaction.
This is the one-sorted standard-part construction of
\cite[Ch.~II, \S2.6, pp.~39--40]{Ebbinghaus1985ExtendedLogics}.
The brackets distinguish a restriction on permitted structures from
the operator extension \(\operatorname{FO}(Q_K^*)\).
\end{hbdefinition}

\begin{hbtheorem}{thm:p03-lq-standard-omega}{proved=annex:p03.lq.6}{Standard parts and projective representation}
\begin{enumerate}
 \item Every sentence of \(\operatorname{FO}[K]\), with its
       standard-part satisfaction condition, has an equivalent
       sentence of \(\operatorname{FO}(Q_K^*)\).
 \item Let \(K_\omega\) be the class of orders isomorphic to
       \((\omega,<)\). For a vocabulary disjoint from the fixed
       standard-part symbols, the infinite models of any
       \(\operatorname{FO}(Q_{K_\omega}^*)\)-sentence form a PC
       class in \(\operatorname{FO}[K_\omega]\).
       Its models of all nonempty cardinalities form an RPC class
       in that standard-part semantics.
 \item Under the same-domain definition of PC in
       \cref{def:p03-w03-p03-047-ec-pc-rpc}, the second clause cannot
       be strengthened to a PC representation of all nonempty models:
       even a tautology's finite models cannot be such reducts.
\end{enumerate}
The finite-domain qualification is essential for the one-sorted
projective comparison discussed in
\cite[Ch.~II, \S4.1, p.~52]{Ebbinghaus1985ExtendedLogics}.
\end{hbtheorem}

\begin{hbtheorem}{thm:p03-lq-field-boundary}{proved=annex:p03.lq.7}{Different compactness behaviours}
Let \(K\) be the class of characteristic-zero fields in the usual
finite field vocabulary. Then \(\operatorname{FO}[K]\) is compact,
whereas \(\operatorname{FO}(Q_K^*)\) is not countably compact.
The failure in the latter has witnesses to finite satisfiability
with infinite domains. Thus restricting models to have a standard
part and adding the corresponding class quantifier are different
operations even at the level of compactness
\cite[Ch.~II, \S4.1, p.~52]{Ebbinghaus1985ExtendedLogics}.
\end{hbtheorem}

\subsection{Monotone quantifiers and finite comparisons}

\begin{hbdefinition}{def:p03-lq-monotone}{Monotone classes and partial maps}
For a finite relational pattern \(\sigma\), a class \(K\) is
\emph{monotone} if enlarging any or all of its relations on a fixed
domain preserves membership in \(K\). The relation interpretations
are enlarged componentwise; the domain is not changed.

A \emph{graph partial isomorphism} \(p\) from \(\mathfrak A\) to
\(\mathfrak B\) is an injective map from a subset of \(A\) into \(B\)
which preserves and reflects the following atoms whenever all their
displayed variable values lie in its domain:
\[
 R(\bar x),\qquad x=y,\qquad c=y,\qquad f(\bar x)=y.
\]
Nullary relation atoms are included. The domain need not be closed
under functions or contain the constants. A formula is
\emph{term-reduced} if these are its only atomic forms.
Its quantifier rank is zero at atoms, unchanged by negation, the
maximum at a finite Boolean combination, and increased by one at
an existential quantifier or a class-quantifier application.
At the latter, use the maximum of the ranks of its operands;
the maximum of an empty list is zero. Empty conjunction and disjunction
are the rank-zero logical constants true and false.
\end{hbdefinition}

\begin{hbdefinition}{def:p03-lq-systems}{Finite and unbounded comparison systems}
Fix a family of monotone quantifiers with finite relational patterns.
An \(n\)-system between \(\mathfrak A\) and \(\mathfrak B\) is a
sequence \(I_0,\ldots,I_n\) of nonempty sets of finite graph partial
isomorphisms, satisfying the following for every \(p\in I_{m+1}\),
where \(m<n\).
\begin{enumerate}
 \item For each \(a\in A\) there is \(q\in I_m\) extending \(p\)
       with \(a\) in its domain; symmetrically every \(b\in B\)
       can be included in the range of such an extension.
 \item For each quantifier \(Q_K\) and each pattern structure
       \(\mathfrak C\in K\) with domain \(A\), there is
       \(\mathfrak D\in K\) with domain \(B\) such that, for every
       pattern relation \(R\) and \(\bar d\in R^{\mathfrak D}\),
       some \(q\in I_m\) extends \(p\), includes all entries of
       \(\bar d\) in its range, and has
       \(q^{-1}(\bar d)\in R^{\mathfrak C}\).
       The symmetric condition, starting with \(\mathfrak D\) on
       \(B\), is also required.
\end{enumerate}
An \(\omega\)-system is a sequence \((I_m)_{m<\omega}\) satisfying
these conditions at every finite level. An \emph{unbounded system}
is a nonempty set \(I\) of finite graph partial isomorphisms with
the same extension conditions, but with every responding map again
in \(I\), without decreasing a level.
The empty tuple in a true nullary relation is treated by the same
rule; no tuple entries then have to be added.
\end{hbdefinition}

\begin{hbtheorem}{thm:p03-lq-finite-comparison}{proved=annex:p03.lq.8}{Finite comparison theorem}
Let the ambient vocabulary be finite and let \(Q_1,\ldots,Q_s\)
be a finite family of monotone quantifiers with finite relational
patterns. For every \(n<\omega\), the following are equivalent:
\begin{enumerate}
 \item \(\mathfrak A\) and \(\mathfrak B\) agree on all term-reduced
       sentences of \(\operatorname{FO}(Q_1,\ldots,Q_s)\) of rank
       at most \(n\).
 \item There is an \(n\)-system between them.
\end{enumerate}
Consequently they agree on all sentences of this logic if and only if
there is an \(n\)-system for every finite \(n\), and if and only if
there is an \(\omega\)-system. Arbitrary terms are allowed in this
last sentence-equivalence assertion: term reduction preserves truth,
but can increase rank.
For no added quantifiers this recovers the first-order comparison
theorem. This is the monotone-quantifier comparison of
\cite[Ch.~II, Theorem~4.2.3 and Corollary~4.2.4,
pp.~54--56]{Ebbinghaus1985ExtendedLogics}.
\end{hbtheorem}

\begin{hbdefinition}{def:p03-lq-characteristic}{The source's characteristic formulas}
Fix a finite ambient vocabulary and one monotone quantifier \(Q_K\)
for a single relation of arity \(r\geq1\). For \(\bar a\in A^k\),
let \(\xi^0_{\bar a}(\bar x)\) be its complete graph-atomic diagram
in the displayed variables: the conjunction of true graph atoms and
negations of false ones. For a set \(M\) of tuples of one fixed
finite length, abbreviate
\[
 O^m_{\bar a,M}(\bar x,\bar y)
       =\bigvee_{\bar c\in M}\xi^m_{\bar a\bar c}(\bar x,\bar y).
\]
For length one, write \(A\) for its set of one-tuples. Define
\[
\begin{aligned}
 \xi^{m+1}_{\bar a}(\bar x)={}&
 \bigwedge_{c\in A}\exists y\,\xi^m_{\bar a c}(\bar x,y)
       \\
 &\land\forall y\,O^m_{\bar a,A}(\bar x,y)\\
 &\land\!\bigwedge_{\substack{M\subseteq A^r\\ (A,M)\in K}}
             Q_K\bar y\,O^m_{\bar a,M}(\bar x,\bar y)\\
 &\land\!\bigwedge_{\substack{M\subseteq A^r\\ (A,M)\notin K}}
             \neg Q_K\bar y\,\neg
                     O^m_{\bar a,A^r\setminus M}(\bar x,\bar y).
\end{aligned}
\]
All displayed indexed conjunctions and disjunctions are replaced by
finite lists of their equivalence-distinct operands. The finite-type
argument in \cref{annex:p03.lq.8} justifies these finite choices;
no infinitary connective is being added. Retain the displayed
quantifier wrappers, even when a formula happens to be a tautology.
This is the construction denoted \(\psi^m\) in
\cite[Ch.~II, Definition~4.2.5, p.~55]{Ebbinghaus1985ExtendedLogics}.
\end{hbdefinition}

\begin{hbtheorem}{thm:p03-lq-characteristic}{proved=annex:p03.lq.14}{Characteristic formulas and types}
The formula \(\xi^m_{\bar a}\) has rank \(m\), is true at
\(\bar a\) in \(\mathfrak A\), and defines the complete
rank-at-most-\(m\) term-reduced type of that tuple.
Consequently
\(\xi^m_{\bar a\bar c}(\bar x,\bar y)\models\xi^m_{\bar a}(\bar x)\)
and
\(\xi^{m+1}_{\bar a}\models\xi^m_{\bar a}\).
Thus these formulas are equivalent to the characteristic formulas
used in \cref{annex:p03.lq.8}, with the source's syntactic rank and
projection properties retained
\cite[Ch.~II, Lemma~4.2.6, p.~56]{Ebbinghaus1985ExtendedLogics}.
\end{hbtheorem}

\subsection{Infinitary comparisons}

\begin{hbdefinition}{def:p03-lq-infinitary}{Infinitary formulas with class quantifiers}
The language \(L_{\infty\omega}(Q_i:i\in J)\) additionally admits
conjunctions and disjunctions indexed by arbitrary sets, provided the
union of their free-variable sets is finite. Quantifier applications
retain finite pattern vocabularies and finite bound-variable lists.
Formula trees are well-founded. The ambient vocabulary and the
family of quantifiers need not be finite.
\end{hbdefinition}

\begin{hbtheorem}{thm:p03-lq-karp}{proved=annex:p03.lq.9}{Infinitary comparison theorem}
For any vocabulary and any family of monotone finite relational
class quantifiers, two structures satisfy the same sentences of
\(L_{\infty\omega}(Q_i:i\in J)\) if and only if there is an
unbounded comparison system between them. This is the generalized
Karp theorem of \cite[Ch.~II, Theorem~4.3.3,
p.~59]{Ebbinghaus1985ExtendedLogics}.
\end{hbtheorem}

\begin{hbcorollary}{thm:p03-lq-countable}{proved=annex:p03.lq.10}{Ordinary comparison systems}
With no added quantifiers, an unbounded graph comparison system is
an ordinary partial isomorphism system in the sense of
\cref{def:p03-w03-p03-039-partial-iso}. In particular, two at most
countable structures linked by such a system are isomorphic.
This countability condition cannot be omitted.
Moreover, if a logic \(\mathscr R\) gives the same sentences to any
two ordinarily partially isomorphic structures, then
\(\mathscr R\leq_{\equiv}L_{\infty\omega}\) in the sense of
\cref{def:p03-w03-p03-047-strength}. This is a comparison of
which structures can be distinguished, not a claimed translation
of every \(\mathscr R\)-sentence into one infinitary sentence.
\end{hbcorollary}

\subsection{Uncountability moves and interpolation}

\begin{hbdefinition}{def:p03-lq-qone-game}{The finite uncountability game}
The game for \(Q_1\) has \(n\) rounds and a current finite map of
selected elements. In each round the first player chooses either:
\begin{enumerate}
 \item an element in either structure; the second player answers
       with an element in the other; or
 \item an uncountable subset of either structure; the second player
       answers with an uncountable subset of the other, the first
       player chooses an element of this answering subset, and the
       second player chooses its mate from the original subset.
\end{enumerate}
The selected pair extends the current map. The second player wins
if the resulting map after all rounds is a graph partial isomorphism;
an inconsistent or noninjective choice loses immediately.
If no answering uncountable subset exists, the second player loses
that move. This is the game described in
\cite[Ch.~II, \S4.2, p.~57]{Ebbinghaus1985ExtendedLogics}.
\end{hbdefinition}

\begin{hbtheorem}{thm:p03-lq-qone-game}{proved=annex:p03.lq.11}{Games and uncountability}
In finite vocabulary, two structures agree on all sentences of
\(\operatorname{FO}(Q_1)\) if and only if the second player has a
winning strategy in the finite uncountability game for every
\(n<\omega\).
\end{hbtheorem}

\begin{hbtheorem}{thm:p03-lq-qone-interpolation}{proved=annex:p03.lq.12}{Failure of interpolation for uncountability}
Let \(\mathfrak A\) and \(\mathfrak B\) be equivalence-relation
structures with infinitely many classes, each class uncountable.
There is an unbounded \(Q_1\)-comparison system between them.
Consequently they agree on every sentence of
\(L_{\infty\omega}(Q_1)\), regardless of the number of their classes.
In particular, the entailment of
\cref{def:p03-w03-p03-047-keisler-entailment} has no interpolant
in \(\operatorname{FO}(Q_1)\).
This supplies the comparison underlying
\cite[Ch.~II, Application~4.2.7, p.~57]{Ebbinghaus1985ExtendedLogics}.
\end{hbtheorem}

\subsection{Ordinary comparison systems in algebra and order}

\begin{hbtheorem}{thm:p03-lq-examples}{proved=annex:p03.lq.13}{Examples of ordinary comparison systems}
\begin{enumerate}
 \item Any two infinite structures of the empty vocabulary have an
       unbounded ordinary partial isomorphism system.
 \item Any two dense linear orders without endpoints have such a
       system. For every uncountable cardinal \(\kappa\), there
       are two nonisomorphic such orders of cardinality \(\kappa\).
 \item Any two algebraically closed fields of the same characteristic,
       each of infinite transcendence degree over its prime field,
       have such a system.
\end{enumerate}
These are the ordinary partial-isomorphism examples in
\cite[Ch.~II, \S4.3, pp.~57--58]{Ebbinghaus1985ExtendedLogics}.
\end{hbtheorem}


\section{Compactness, vocabulary and spaces of theories}
\label{sec:p03-compactness-vocabulary}

Compactness bounds the information needed to distinguish structures.
It also has a topological formulation, but the fixed-cardinal versions
must distinguish single sentences from whole theories placed in blocks.

\begin{hbdefinition}{def:p03-cv-cardinals}{Two cardinal parameters}
Let \(\kappa\geq\lambda\geq\aleph_0\). A logic is
\emph{\((\kappa,\lambda)\)-compact} if, for every vocabulary and
every theory \(T\) with \(|T|\leq\kappa\), satisfiability of
every subset of \(T\) of size less than \(\lambda\) implies
satisfiability of \(T\). Thus \((\kappa,\aleph_0)\)-compactness
is the \(\kappa\)-compactness of
\cref{def:p03-w03-p03-041-compactness}.
The cardinality bound on the whole theory is non-strict; the bound
on its tested subtheories is strict
\cite[Ch.~II, Definition~5.1.1, p.~59]{Ebbinghaus1985ExtendedLogics}.
\end{hbdefinition}

\subsection{The essential vocabulary of a sentence}

\begin{hbdefinition}{def:p03-cv-support}{Semantic support}
A sentence \(\varphi\in\mathscr L[\tau]\) \emph{depends only on}
a sub-vocabulary \(\sigma\subseteq\tau\) if any two
\(\tau\)-structures whose \(\sigma\)-reducts are isomorphic agree
on \(\varphi\). Such a \(\sigma\) is a semantic support.
This says something about satisfaction; it does not assert that the
same written sentence belongs to \(\mathscr L[\sigma]\).
Sort symbols and the sorts attached to retained symbols are included
when counting a vocabulary.
\end{hbdefinition}

\begin{hbtheorem}{thm:p03-cv-support}{proved=annex:p03.cv.1}{Compactness bounds semantic support}
Let \(\mathscr L\) satisfy the general framework of
\cref{def:p03-w03-p03-047-ebbinghaus-logic} and the atom, Boolean
and particularization properties of
\cref{def:p03-w03-p03-047-closure}. It may be many-sorted.
If it is \((\kappa,\lambda)\)-compact and \(|\tau|\leq\kappa\),
every sentence in \(\mathscr L[\tau]\) has a semantic support
of cardinality less than \(\lambda\).
In particular every sentence of a compact logic with these properties
has finite semantic support. This is the support conclusion of
\cite[Ch.~II, Proposition~5.1.2, p.~60]{Ebbinghaus1985ExtendedLogics},
with the sort identifications made explicit.
\end{hbtheorem}

\begin{hbcorollary}{thm:p03-cv-support-translation}{proved=annex:p03.cv.2}{A sentence over its semantic support}
Let \(\mathscr L\) be one-sorted and regular. If
\(\varphi\in\mathscr L[\tau]\) has semantic support
\(\sigma\subseteq\tau\), then some
\(\psi\in\mathscr L[\sigma]\) is equivalent to \(\varphi\)
over all \(\tau\)-structures. Consequently a compact one-sorted
regular logic has an equivalent finite-vocabulary sentence for every
sentence, even when its syntax does not require finite occurrence.
\end{hbcorollary}

\begin{hbexample}{def:p03-cv-unbounded-support}{No uniform support bound for the infinitary language}
For an arbitrary infinite set \(J\), take unary predicates \(P_j\)
and the \(L_{\infty\omega}\)-sentence
\[
 \bigwedge_{j\in J}\forall x\,P_j(x).
\]
Every semantic support must contain every \(P_j\): on a one-element
domain, changing just an omitted predicate from full to empty changes
the sentence's truth without changing the supported reduct.
As \(|J|\) is arbitrary, no single cardinal bounds the essential
vocabularies of all sentences of this language.
\end{hbexample}

\subsection{Spaces of realized complete theories}

\begin{hbdefinition}{def:p03-cv-space}{The theory space of a vocabulary}
Fix \(\tau\) such that \(\mathscr L[\tau]\) is a set, and assume
Boolean closure and first-order tautologies. Define
\[
 \begin{aligned}
 X_{\mathscr L,\tau}
   &=\{\operatorname{Th}_{\mathscr L}(\mathfrak A):
                    \mathfrak A\in\operatorname{Str}[\tau]\},\\
 U_\varphi&=\{p\in X_{\mathscr L,\tau}:\varphi\in p\}.
 \end{aligned}
\]
This is a set: it is a subclass of the power set of
\(\mathscr L[\tau]\). Give it the topology generated by the
\(U_\varphi\). Its points are complete theories actually satisfied
by structures, not arbitrary unfulfilled truth assignments.
Equivalently, choose one structure from each logical-equivalence
class and use the sentence-defined model classes as basic opens.
The two presentations identify corresponding points and basic opens
\cite[Ch.~II, \S5.1, p.~59]{Ebbinghaus1985ExtendedLogics}.
\end{hbdefinition}

\begin{hbtheorem}{thm:p03-cv-space}{proved=annex:p03.cv.3}{Logical and topological compactness}
The sets \(U_\varphi\) form a clopen basis and
\(X_{\mathscr L,\tau}\) is Hausdorff. It is compact if and only if
every finitely satisfiable theory in vocabulary \(\tau\) is
satisfiable. Consequently, for a small logic (each
\(\mathscr L[\tau]\) a set), full logical compactness is equivalent
to compactness of every such theory space.
This extends the first-order Stone-space argument used in
\cref{thm:p03-robinson-joint-consistency}; first-order type-space
compactness itself is \cref{thm:type-space-compact}.
\end{hbtheorem}

\begin{hbdefinition}{def:p03-cv-blocks}{Topological and block compactness}
A space is \emph{\((\kappa,\lambda)\)-compact} if every family
of at most \(\kappa\) closed sets with empty intersection has a
subfamily of size less than \(\lambda\) with empty intersection.
For the fixed vocabulary \(\tau\), the corresponding
\emph{block condition} says: for every family \((T_i)_{i\in I}\)
of theories with \(|I|\leq\kappa\), if
\(\bigcup_{i\in J}T_i\) is satisfiable for each \(J\subseteq I\)
of size less than \(\lambda\), then
\(\bigcup_{i\in I}T_i\) is satisfiable. Each block \(T_i\) is an
arbitrary subset of \(\mathscr L[\tau]\); there is no bound of
\(\kappa\) on the size of a block or of their union.
\end{hbdefinition}

\begin{hbtheorem}{thm:p03-cv-blocks}{proved=annex:p03.cv.4}{Closed sets and theory blocks}
The theory space \(X_{\mathscr L,\tau}\) is topologically
\((\kappa,\lambda)\)-compact exactly when the block condition
holds. Ordinary logical \((\kappa,\lambda)\)-compactness tests
single sentences, equivalently families of sentence-defined basic
clopens; it need not imply this condition for arbitrary closed sets.
If \(|\mathscr L[\tau]|\leq\kappa\), however, logical
\((\kappa,\lambda)\)-compactness does imply the block condition
at this vocabulary.
\end{hbtheorem}

\begin{hbtheorem}{thm:p03-cv-qone-space}{proved=annex:p03.cv.5}{A compactness boundary}
The logic \(\operatorname{FO}(Q_1)\) is countably compact, but
for the vocabulary \(\{c_\alpha:\alpha<\omega_1\}\) its space
of realized complete theories is not topologically
\((\aleph_0,\aleph_0)\)-compact. The latter failure is witnessed
by countably many closed sets whose every finite subfamily even
has a point represented by a finite structure.
No continuum hypothesis is assumed.
\end{hbtheorem}

Thus the unrestricted correspondence in \cref{thm:p03-cv-space}
does not transfer unchanged to fixed cardinal bounds. The extra
strength lies in treating an entire theory as one closed-set block,
as in the strengthened compactness discussed in
\cite[Ch.~II, \S5.1, p.~60]{Ebbinghaus1985ExtendedLogics}.

\section{Ordinal operations and the well-ordering threshold}
\label{sec:p03-ordinal-pinning}

Work with the one-sorted regular logics, nonempty structures and
arbitrary set-sized vocabularies of
\cref{def:p03-general-pinning-setting}. Let \(\kappa\geq1\)
be a cardinal. A theory has at most \(\kappa\) sentences; the
vocabulary is not subject to this bound.

\begin{hbdefinition}{def:p03-op-carrier}{Carrier pinning}
A theory \(T\), with a designated unary predicate \(P\) and binary
relation \(<\), \emph{pins an ordinal on a carrier} if every model
of \(T\) makes \(<\) restricted to \(P\) a strict well-order,
and some model gives that carrier the specified order type.
The carrier may be empty even though the ambient structure is not.
Write \(W_\kappa\) for the supremum of these ordinals, or
\(\infty\) when they are unbounded in the ordinals.
\end{hbdefinition}

\begin{hbtheorem}{thm:p03-op-initial}{proved=annex:p03.op.1}{Initial segments, successors and carriers}
Carrier-pinnable ordinals are closed downwards and under successor.
Every finite ordinal is carrier-pinnable with one sentence.
Moreover \(W_\kappa=w_\kappa(\mathscr L)\), where the latter
uses the field convention of \cref{def:p03-w03-p03-041-pinning}.
If this common value is an ordinal \(W\), it is an infinite limit,
is not itself pinnable, and the carrier-pinnable ordinals are
exactly those below \(W\). Under the strict-field convention,
exactly the same ordinals are pinnable except for \(1\).
\end{hbtheorem}

\begin{hbtheorem}{thm:p03-op-threshold}{proved=annex:p03.op.2}{The cofinal-model threshold}
If \(W_\kappa=W\) is an ordinal, then \(W\) is the least ordinal
\(\gamma\) with the following property: whenever a theory of at
most \(\kappa\) sentences has well-ordered carrier models whose
order types below \(\gamma\) are cofinal in \(\gamma\), it also
has a model whose designated carrier order is not well-ordered.
If \(W_\kappa=\infty\), no ordinal has this property.
For strict fields rather than designated carriers, \(W\) is the
least \emph{infinite} ordinal with the corresponding property.
\end{hbtheorem}

The qualification about fields is necessary: the field of a strict
order cannot have exactly one element. Consequently the condition
on types cofinal below \(2\) is impossible for fields, and the
unqualified least-ordinal characterization would be vacuous there.
Carrier pinning gives the full threshold formulation without
changing the well-ordering number. Compare
\cite[Ch.~II, Proposition~5.2.2, pp.~61--62]{Ebbinghaus1985ExtendedLogics}.

\begin{hbtheorem}{thm:p03-op-arithmetic}{proved=annex:p03.op.3}{Ordinal sum, product and exponentiation}
If \(\alpha\) and \(\beta\) are carrier-pinnable by theories of
at most \(\kappa\) sentences, so are
\(\alpha+\beta\), \(\alpha\beta\), and \(\alpha^\beta\),
with \(0^0=1\). In particular, a bounded \(W_\kappa\) is
closed under these operations on its smaller ordinals.
\end{hbtheorem}


\section{Interpolation and the Robinson property}
\label{sec:p03-w03-p03-047-robinson}

Ebbinghaus's last section works with regular logics, one-sorted and
many-sorted, and writes (R)PC for PC in the one-sorted and RPC in the
many-sorted case \cite[Ch.~II, \S7, p.~68]{Ebbinghaus1985ExtendedLogics}.

\begin{hbdefinition}{def:p03-w03-p03-047-interpolation-for}{Interpolation for another logic (Ebbinghaus, Definition~7.1.1)}
For logics \(\mathscr L\) and \(\mathscr L^\ast\), \(\mathscr L^\ast\)
\emph{has the interpolation property for} \(\mathscr L\) (\emph{allows
interpolation for} \(\mathscr L\)) if any two disjoint classes of
structures of one vocabulary that are (R)PC in \(\mathscr L\)
(\cref{def:p03-w03-p03-047-ec-pc-rpc}) are separated by a class that is EC
in \(\mathscr L^\ast\) \cite[Ch.~II, Definition~7.1.1, p.~68]{Ebbinghaus1985ExtendedLogics}.
\end{hbdefinition}

For \(\mathscr L^\ast=\mathscr L\) one-sorted this is the separation form of
interpolation in \cref{thm:p03-w03-p03-047-craig-separation}.  Ebbinghaus's
positive examples are first-order logic (Craig in the one-sorted case,
Feferman in the many-sorted case) and \(\mathscr L_{\omega_1\omega}\) with
its countable admissible fragments (Lopez-Escobar, Barwise)
\cite[Ch.~II, Examples~7.1.2, p.~68]{Ebbinghaus1985ExtendedLogics}; for the
one-sorted cases see \cref{thm:p02-w02c-p02-003-interpolation}.

\begin{hbdefinition}{def:p03-w03-p03-047-robinson-property}{Robinson property (Ebbinghaus, Definition~7.1.4)}
\(\mathscr L\) has the \emph{Robinson property} if for all vocabularies
\(\tau_0,\tau_1\) with \(\tau=\tau_0\cap\tau_1\), and all classes
\(\Phi\subseteq\mathscr L[\tau]\), \(\Phi_0\subseteq\mathscr L[\tau_0]\),
\(\Phi_1\subseteq\mathscr L[\tau_1]\): if \(\Phi\) is \emph{complete} (all
\(\tau\)-models of \(\Phi\) are \(\mathscr L\)-equivalent) and
\(\Phi\cup\Phi_i\) has a model for \(i=0,1\), then
\(\Phi\cup\Phi_0\cup\Phi_1\) has a model
\cite[Ch.~II, Definition~7.1.4, p.~70]{Ebbinghaus1985ExtendedLogics}.
\(\mathscr L\) is \emph{small} if every \(\mathscr L[\tau]\) is a set.
\end{hbdefinition}

\begin{hbtheorem}{thm:p03-w03-p03-047-robinson-interpolation}{proved=annex:p03-w03-p03-047-robinson-interpolation}{Robinson property and interpolation for compact logics (Ebbinghaus, Proposition~7.1.5)}
Let \(\mathscr L\) be a small, compact logic with the atom, Boolean and
particularization properties (\cref{def:p03-w03-p03-047-closure}), for
instance a small, compact, regular logic.  Then \(\mathscr L\)
has the interpolation property
(\cref{def:p03-w03-p03-047-ebbinghaus-properties}~(viii)) iff it has the
Robinson property \cite[Ch.~II, Proposition~7.1.5, p.~70]{Ebbinghaus1985ExtendedLogics}.
\end{hbtheorem}

Ebbinghaus adds that for logics with a sufficiently small occurrence number
the Robinson property implies compactness, a result of later chapters of
the volume \cite[Ch.~II, \S7.1, p.~70]{Ebbinghaus1985ExtendedLogics}; with it,
Barwise's formulation follows: for logics whose sentences contain few
symbols, the Robinson property is equivalent to compactness together with
interpolation \cite[Ch.~I, \S2.4, p.~16]{Barwise1985BackgroundAims}.  That
implication from the Robinson property to compactness is not proved here.
Without compactness the two properties separate, as Barwise reports for
\(\mathscr L_{\omega_1\omega}\) \cite[Ch.~I, \S2.3, p.~16]{Barwise1985BackgroundAims};
the next theorem proves this together with Ebbinghaus's first
counterexample to interpolation between logics.  Write
\(\mathscr L_{\infty\omega}\) for the union of the logics
\(\mathscr L_{\kappa\omega}\) of \cref{def:p03-w03-p03-047-lkl} over all
infinite cardinals \(\kappa\).

\begin{hbtheorem}{thm:p03-w03-p03-047-infinite-sets}{proved=annex:p03-w03-p03-047-infinite-sets}{Infinite sets in \(\mathscr L_{\infty\omega}\); Robinson fails for \(\mathscr L_{\omega_1\omega}\)}
\begin{enumerate}
 \item Any two infinite sets, as structures of the empty vocabulary,
       satisfy the same \(\mathscr L_{\infty\omega}\)-sentences.
 \item \(\mathscr L_{\infty\omega}\) does not allow interpolation for
       \(\mathscr L_{\omega_2\omega}\): the classes
       \(\mathfrak K_{\aleph_0}\) of nonempty sets of power at most
       \(\aleph_0\) and \(\mathfrak K^{\aleph_1}\) of sets of power at least
       \(\aleph_1\) are disjoint and PC in \(\mathscr L_{\omega_2\omega}\), and
       no class EC in \(\mathscr L_{\infty\omega}\) separates them
       \cite[Ch.~II, Counterexamples~7.1.3(a), p.~69]{Ebbinghaus1985ExtendedLogics}.
 \item \(\mathscr L_{\omega_1\omega}\) does not have the Robinson property,
       although it has the interpolation property
       (\cref{thm:p02-w02c-p02-003-interpolation}).
\end{enumerate}
\end{hbtheorem}

Part~3 also shows that \(\mathscr L_{\omega_1\omega}\) is not compact, by
\cref{thm:p03-w03-p03-047-robinson-interpolation} read backwards:
\(\mathscr L_{\omega_1\omega}\) is small and has these closure properties, so if
it were compact its interpolation property would give it the Robinson
property.
Ebbinghaus's argument for part~2 cites Karp's theorem on partially
isomorphic structures; the proof given here verifies the needed instance
directly by induction on formulas.
